\section*{Lecture 2: Information and Decision Modeling}
\addcontentsline{toc}{section}{Lecture 2: Information and Decision Modeling}
\clearpage
\SlideHeader{Lecture 2: Information and Decision Modeling}{001}{表紙}
\SlideImage{1}{../Materials/2026/3DM_02_Information-and-Decision-Modeling_SS26.pdf}
\begin{multicols}{2}\scriptsize\raggedcolumns
\NoteSection{日本語訳}
\begin{itemize}
\item Data Driven Decision Making
\item 第2回 - 情報モデリングと意思決定モデリング
\end{itemize}
\end{multicols}

\clearpage
\SlideHeader{Lecture 2: Information and Decision Modeling}{002}{本日の講義目標}
\SlideImage{2}{../Materials/2026/3DM_02_Information-and-Decision-Modeling_SS26.pdf}
\begin{multicols}{2}\footnotesize\raggedcolumns
\NoteSection{日本語訳}
\begin{itemize}
\item Business Analytics のプロセス。
\item 情報モデリングと意思決定モデリング。
\item 時間依存旅行時間を持つ車両ルーティング。
\item データが最初に来る。
\item 注意: この講義では多くの手法を使うが、詳細な定義ではなく考え方のみを扱う。
\end{itemize}
\NoteSection{解説}
Real-World は、実際の配送現場、道路、顧客、車両、作業員のような現実対象である。現実は非常に複雑なので、そのまま最適化モデルには入れられない。\par

Aggregation は、現実データを意思決定で使える形に圧縮する処理である。例: GPS点列を道路セグメント別速度へ変換する、住所を顧客ノードへ変換する、過去走行記録から旅行時間行列を作る。\par

Information Model は『意思決定に必要な情報の表現』である。顧客集合、デポ、道路ネットワーク、旅行時間行列、需要分布、時間帯などが入る。Decision Model は『その情報を使って何を最適化するか』を表し、決定変数、目的関数、制約、実行可能解を定義する。\par
\NoteSection{試験対策ポイント}
\begin{itemize}
\item Real-World -> Aggregation -> Information Model -> Decision Model -> Implementation -> Real-World の流れを、配送例で説明できる。
\item Information Model と Decision Model を混同しない。前者は情報表現、後者は意思決定の評価・選択構造である。
\end{itemize}
\NoteSection{確認問題と解答}
\begin{enumerate}
\item \textbf{Information Model と Decision Model の違いを説明せよ。}\\Information Model は、現実を意思決定に必要な形へ圧縮した情報表現である。例として、顧客集合、距離行列、時間帯別旅行時間、需要分布がある。Decision Model は、その情報を用いて選択肢を評価する数学的・論理的構造であり、決定変数、目的関数、制約を持つ。配送例なら、旅行時間行列は情報モデル、訪問順序を選んで総旅行時間を最小化するTSP/VRPは意思決定モデルである。
\item \textbf{Aggregation と Implementation を配送例で説明せよ。}\\Aggregation は、GPS記録、道路地図、住所データなどを、顧客ノードや旅行時間行列に変換する処理である。Implementation は、最適化モデルが出した訪問順序やアーク選択を、実際のドライバーが使える配送順、住所リスト、ナビゲーションに変換する処理である。Aggregation は現実からモデルへ、Implementation はモデルから現実への橋である。
\end{enumerate}
\end{multicols}

\clearpage
\SlideHeader{Lecture 2: Information and Decision Modeling}{003}{アジェンダ}
\SlideImage{3}{../Materials/2026/3DM_02_Information-and-Decision-Modeling_SS26.pdf}
\begin{multicols}{2}\scriptsize\raggedcolumns
\NoteSection{日本語訳}
\begin{itemize}
\item 情報モデリングと意思決定モデリング。
\item ケーススタディ: 時間依存旅行時間を持つ巡回セールスマン問題。
\end{itemize}
\end{multicols}

\clearpage
\SlideHeader{Lecture 2: Information and Decision Modeling}{004}{意思決定}
\SlideImage{4}{../Materials/2026/3DM_02_Information-and-Decision-Modeling_SS26.pdf}
\begin{multicols}{2}\footnotesize\raggedcolumns
\NoteSection{日本語訳}
\begin{itemize}
\item Real-World: 行動が実行され、その結果として状態を観測できる。
\item Aggregation: 現実世界のデータを情報モデルの次元へ移す。
\item Information Model: 情報を圧縮された形で描写する。
\item Decision Model: 決定空間を定義し、決定を評価する。
\item Implementation: 意思決定モデルの解を現実世界の行動へ変換する。
\end{itemize}
\NoteSection{解説}
Real-World は、実際の配送現場、道路、顧客、車両、作業員のような現実対象である。現実は非常に複雑なので、そのまま最適化モデルには入れられない。\par

Aggregation は、現実データを意思決定で使える形に圧縮する処理である。例: GPS点列を道路セグメント別速度へ変換する、住所を顧客ノードへ変換する、過去走行記録から旅行時間行列を作る。\par

Information Model は『意思決定に必要な情報の表現』である。顧客集合、デポ、道路ネットワーク、旅行時間行列、需要分布、時間帯などが入る。Decision Model は『その情報を使って何を最適化するか』を表し、決定変数、目的関数、制約、実行可能解を定義する。\par
\NoteSection{試験対策ポイント}
\begin{itemize}
\item Real-World -> Aggregation -> Information Model -> Decision Model -> Implementation -> Real-World の流れを、配送例で説明できる。
\item Information Model と Decision Model を混同しない。前者は情報表現、後者は意思決定の評価・選択構造である。
\end{itemize}
\NoteSection{確認問題と解答}
\begin{enumerate}
\item \textbf{Information Model と Decision Model の違いを説明せよ。}\\Information Model は、現実を意思決定に必要な形へ圧縮した情報表現である。例として、顧客集合、距離行列、時間帯別旅行時間、需要分布がある。Decision Model は、その情報を用いて選択肢を評価する数学的・論理的構造であり、決定変数、目的関数、制約を持つ。配送例なら、旅行時間行列は情報モデル、訪問順序を選んで総旅行時間を最小化するTSP/VRPは意思決定モデルである。
\item \textbf{Aggregation と Implementation を配送例で説明せよ。}\\Aggregation は、GPS記録、道路地図、住所データなどを、顧客ノードや旅行時間行列に変換する処理である。Implementation は、最適化モデルが出した訪問順序やアーク選択を、実際のドライバーが使える配送順、住所リスト、ナビゲーションに変換する処理である。Aggregation は現実からモデルへ、Implementation はモデルから現実への橋である。
\end{enumerate}
\end{multicols}

\clearpage
\SlideHeader{Lecture 2: Information and Decision Modeling}{005}{情報モデルと意思決定モデルの選択}
\SlideImage{5}{../Materials/2026/3DM_02_Information-and-Decision-Modeling_SS26.pdf}
\begin{multicols}{2}\footnotesize\raggedcolumns
\NoteSection{日本語訳}
\begin{itemize}
\item 情報モデルと意思決定モデルは初期仮説に基づいて設計される。
\item 仮説は現実問題の観測行動とデータに基づく。これは descriptive analytics に対応する。
\item 仮説は必ずしも正しいとは限らない。
\item 仮説の品質は、意思決定モデルの解を実装することで評価できる。
\item 観測結果により仮説が更新され、その結果として情報モデルと意思決定モデルも更新される。
\end{itemize}
\NoteSection{解説}
Real-World は、実際の配送現場、道路、顧客、車両、作業員のような現実対象である。現実は非常に複雑なので、そのまま最適化モデルには入れられない。\par

Aggregation は、現実データを意思決定で使える形に圧縮する処理である。例: GPS点列を道路セグメント別速度へ変換する、住所を顧客ノードへ変換する、過去走行記録から旅行時間行列を作る。\par

Information Model は『意思決定に必要な情報の表現』である。顧客集合、デポ、道路ネットワーク、旅行時間行列、需要分布、時間帯などが入る。Decision Model は『その情報を使って何を最適化するか』を表し、決定変数、目的関数、制約、実行可能解を定義する。\par
\NoteSection{試験対策ポイント}
\begin{itemize}
\item Real-World -> Aggregation -> Information Model -> Decision Model -> Implementation -> Real-World の流れを、配送例で説明できる。
\item Information Model と Decision Model を混同しない。前者は情報表現、後者は意思決定の評価・選択構造である。
\end{itemize}
\NoteSection{確認問題と解答}
\begin{enumerate}
\item \textbf{Information Model と Decision Model の違いを説明せよ。}\\Information Model は、現実を意思決定に必要な形へ圧縮した情報表現である。例として、顧客集合、距離行列、時間帯別旅行時間、需要分布がある。Decision Model は、その情報を用いて選択肢を評価する数学的・論理的構造であり、決定変数、目的関数、制約を持つ。配送例なら、旅行時間行列は情報モデル、訪問順序を選んで総旅行時間を最小化するTSP/VRPは意思決定モデルである。
\item \textbf{Aggregation と Implementation を配送例で説明せよ。}\\Aggregation は、GPS記録、道路地図、住所データなどを、顧客ノードや旅行時間行列に変換する処理である。Implementation は、最適化モデルが出した訪問順序やアーク選択を、実際のドライバーが使える配送順、住所リスト、ナビゲーションに変換する処理である。Aggregation は現実からモデルへ、Implementation はモデルから現実への橋である。
\end{enumerate}
\end{multicols}

\clearpage
\SlideHeader{Lecture 2: Information and Decision Modeling}{006}{IDA: Appearance から Structureへ}
\SlideImage{6}{../Materials/2026/3DM_02_Information-and-Decision-Modeling_SS26.pdf}
\begin{multicols}{2}\footnotesize\raggedcolumns
\NoteSection{日本語訳}
\begin{itemize}
\item 現実システムからのデータは、その見かけ appearance を反映するだけで、構造は直接見えにくい。
\item appearance は不完全、誤り、例外を含む可能性がある。
\item システム構造はデータモデル構築を助ける。
\item 記録データは一般的なシステムappearanceの少数サンプルにすぎない。
\item 記録データは、クリーニング、修正、補完、検証、絞り込みによりターゲットデータセットへ変換される。
\item データマイニングはターゲットデータからモデルを計算する。
\item データマイニング概念は、データ内の統計的相関に関する仮説を実装する。
\item モデルまたはパターンはシステム構造に関する証拠を与える。
\end{itemize}
\NoteSection{解説}
IDAは appearance から structure へ向かう。つまり、記録データの表面的なパターンから、背後の情報構造を推測する。例: GPS走行ログから道路混雑パターンを推定する。\par

PMT/ORは structure から appearance/solution へ向かう。つまり、問題構造を仮定し、目的関数・制約・決定変数を作り、解を求める。例: 顧客集合と旅行時間行列からTSPを作る。\par

Business Analytics では、データ側と意思決定側をつなぐ。データ分析だけでは『何を実行すべきか』が出ない。最適化だけでも、現実のデータを正しく表現しなければ役に立たない。\par
\NoteSection{試験対策ポイント}
\begin{itemize}
\item IDA = データから構造を発見する方向、PMT/OR = 構造から意思決定モデルと解を作る方向。
\item 三分野を列挙するだけでなく、交差領域がなぜ必要かを説明する。
\end{itemize}
\NoteSection{確認問題と解答}
\begin{enumerate}
\item \textbf{IDA と PMT/OR の方向性の違いを説明せよ。}\\IDAは記録データ、つまり appearance から出発し、クリーニング、補完、検証、データマイニングを通じて情報構造やパターンを見つける。PMT/ORは問題構造を仮説として置き、目的、制約、決定空間を定義して、実行可能な解を計算する。DDDMでは両方が必要で、情報モデルが両者を接続する。
\item \textbf{Business Analytics の三分野を使って、配送計画を説明せよ。}\\Statistics は旅行時間や需要の不確実性を推定する。Computer Science/Information Systems はGPS、注文、地図データを保存・処理し、システムとして実装する。Operations Research は訪問順序、車両割当、制約を最適化する。三つを統合することで、データに基づく実行可能な配送計画が得られる。
\end{enumerate}
\end{multicols}

\clearpage
\SlideHeader{Lecture 2: Information and Decision Modeling}{007}{PMT: Structure から Appearanceへ}
\SlideImage{7}{../Materials/2026/3DM_02_Information-and-Decision-Modeling_SS26.pdf}
\begin{multicols}{2}\footnotesize\raggedcolumns
\NoteSection{日本語訳}
\begin{itemize}
\item 現実の意思決定問題は、数理モデリングの要件を満たすよう単純化される。
\item 仮説により重要な問題構造の特徴が選択される。
\item データが利用可能な特徴だけが考慮される。
\item 重要な特徴は目的、決定空間、制約である。
\item 意思決定モデルの一般構造が、問題から意思決定モデルへの道を橋渡しする。
\item 意思決定モデル構造の例は、割当、ルーティング、スケジューリングである。
\item 得られた解は問題の一つの最適な appearance を描く。
\end{itemize}
\NoteSection{解説}
IDAは appearance から structure へ向かう。つまり、記録データの表面的なパターンから、背後の情報構造を推測する。例: GPS走行ログから道路混雑パターンを推定する。\par

PMT/ORは structure から appearance/solution へ向かう。つまり、問題構造を仮定し、目的関数・制約・決定変数を作り、解を求める。例: 顧客集合と旅行時間行列からTSPを作る。\par

Business Analytics では、データ側と意思決定側をつなぐ。データ分析だけでは『何を実行すべきか』が出ない。最適化だけでも、現実のデータを正しく表現しなければ役に立たない。\par
\NoteSection{試験対策ポイント}
\begin{itemize}
\item IDA = データから構造を発見する方向、PMT/OR = 構造から意思決定モデルと解を作る方向。
\item 三分野を列挙するだけでなく、交差領域がなぜ必要かを説明する。
\end{itemize}
\NoteSection{確認問題と解答}
\begin{enumerate}
\item \textbf{IDA と PMT/OR の方向性の違いを説明せよ。}\\IDAは記録データ、つまり appearance から出発し、クリーニング、補完、検証、データマイニングを通じて情報構造やパターンを見つける。PMT/ORは問題構造を仮説として置き、目的、制約、決定空間を定義して、実行可能な解を計算する。DDDMでは両方が必要で、情報モデルが両者を接続する。
\item \textbf{Business Analytics の三分野を使って、配送計画を説明せよ。}\\Statistics は旅行時間や需要の不確実性を推定する。Computer Science/Information Systems はGPS、注文、地図データを保存・処理し、システムとして実装する。Operations Research は訪問順序、車両割当、制約を最適化する。三つを統合することで、データに基づく実行可能な配送計画が得られる。
\end{enumerate}
\end{multicols}

\clearpage
\SlideHeader{Lecture 2: Information and Decision Modeling}{008}{IDAとPMTの統一的見方}
\SlideImage{8}{../Materials/2026/3DM_02_Information-and-Decision-Modeling_SS26.pdf}
\begin{multicols}{2}\footnotesize\raggedcolumns
\NoteSection{日本語訳}
\begin{itemize}
\item 記録データからターゲットデータへ。
\item ターゲットデータから情報構造と情報モデル/パターンへ。
\item 問題構造から意思決定モデル構造と意思決定モデルへ。
\item 仮説が、データ側と意思決定側の両方を接続する。
\item 解は現実で実装され、再び観測につながる。
\end{itemize}
\NoteSection{解説}
IDAは appearance から structure へ向かう。つまり、記録データの表面的なパターンから、背後の情報構造を推測する。例: GPS走行ログから道路混雑パターンを推定する。\par

PMT/ORは structure から appearance/solution へ向かう。つまり、問題構造を仮定し、目的関数・制約・決定変数を作り、解を求める。例: 顧客集合と旅行時間行列からTSPを作る。\par

Business Analytics では、データ側と意思決定側をつなぐ。データ分析だけでは『何を実行すべきか』が出ない。最適化だけでも、現実のデータを正しく表現しなければ役に立たない。\par
\NoteSection{試験対策ポイント}
\begin{itemize}
\item IDA = データから構造を発見する方向、PMT/OR = 構造から意思決定モデルと解を作る方向。
\item 三分野を列挙するだけでなく、交差領域がなぜ必要かを説明する。
\end{itemize}
\NoteSection{確認問題と解答}
\begin{enumerate}
\item \textbf{IDA と PMT/OR の方向性の違いを説明せよ。}\\IDAは記録データ、つまり appearance から出発し、クリーニング、補完、検証、データマイニングを通じて情報構造やパターンを見つける。PMT/ORは問題構造を仮説として置き、目的、制約、決定空間を定義して、実行可能な解を計算する。DDDMでは両方が必要で、情報モデルが両者を接続する。
\item \textbf{Business Analytics の三分野を使って、配送計画を説明せよ。}\\Statistics は旅行時間や需要の不確実性を推定する。Computer Science/Information Systems はGPS、注文、地図データを保存・処理し、システムとして実装する。Operations Research は訪問順序、車両割当、制約を最適化する。三つを統合することで、データに基づく実行可能な配送計画が得られる。
\end{enumerate}
\end{multicols}

\clearpage
\ExamHeader{情報モデル・意思決定モデル・Business Analyticsプロセス}
\ExamQuestion{SS24 Aufgabe 1a [8 Punkte]}
\ExamSubsection{問題内容}
Information Model, Decision Model, Aggregation, Disaggregation の機能を説明する問題。\par
\ExamSubsection{満点答案}
Aggregation は、現実世界から得られる詳細で雑多なデータを、情報モデルで扱える抽象度へ変換する機能である。例えば、GPS点列を道路セグメントの速度に変換し、住所を顧客ノードに変換し、過去走行記録を時間帯別旅行時間に集約する。\par
Information Model は、意思決定に必要な情報を構造化した表現である。現実のすべてを含むのではなく、目的に関連する顧客、デポ、道路、旅行時間、需要、時間帯などを圧縮して表す。Decision Model は、情報モデルを入力として、どの決定を選べるか、どの制約を守るか、何を最小化または最大化するかを定義する。TSPなら、アーク選択変数、訪問制約、部分巡回除去制約、総旅行時間最小化が該当する。\par
Disaggregation は、モデルの解を現実の実行可能な指示へ戻す機能である。例えば、モデル上のノード番号列を住所・道路ルート・ドライバー指示に変換する。したがって、Aggregation は現実からモデルへ、Disaggregation はモデルから現実への橋である。\par
\ExamSubsection{具体例による解説}
顧客A,B,Cへの配送では、住所と地図から旅行時間行列を作るのがAggregation、旅行時間行列と顧客集合がInformation Model、最短訪問順を求めるTSPがDecision Model、A->C->Bという解を実際のナビにするのがDisaggregationである。\par
\ExamSubsection{採点で落とさない要素}
\begin{itemize}
\item 四概念をそれぞれ定義
\item モデル入力・出力の方向を明示
\item 配送/TSP例で対応付け
\item 現実とモデルのループを説明
\end{itemize}
\vspace{2mm}\hrule\vspace{2mm}
\ExamQuestion{WS23/24 Aufgabe 3 [10 Punkte]}
\ExamSubsection{問題内容}
配送車両フリートの業務的ツアープランニングを例に、Real-World, Aggregation, Information Model, Decision Model, Implementation の構成要素とデータフローを説明する問題。\par
\ExamSubsection{満点答案}
Real-World は、配送センター、顧客住所、道路ネットワーク、配送車両、ドライバー、交通状況である。ここで実際の配送が行われ、走行時間や遅延などが観測される。Aggregation では、顧客住所をノードにし、道路地図やGPSから顧客間旅行時間を計算し、必要なら時間帯別旅行時間行列を作る。\par
Information Model は、デポ、顧客集合、車両集合、顧客需要、旅行時間または距離行列、時間帯などから成る。Decision Model は、どの車両がどの顧客をどの順序で訪問するかを決めるVRP/TSPである。決定変数はアーク使用や車両割当を表し、制約は各顧客を一度訪問すること、車両容量、勤務時間、時間窓などである。目的関数は総旅行時間、車両数、遅延、コストの最小化である。\par
Implementation は、得られたルートをドライバーの配送順序、ナビゲーション、出発時刻、作業指示に変換することである。実装後の走行時間や遅延は再び観測され、旅行時間モデルや制約仮説の改善に使われる。\par
\ExamSubsection{具体例による解説}
例えば、モデルが『車両1: Depot-A-C-Depot、車両2: Depot-B-D-Depot』を出した場合、Implementationでは実住所、積載順、ナビ経路、時間窓をドライバーに渡す。もしA-C間が常に遅れるなら、AggregationやInformation Modelを更新する。\par
\ExamSubsection{採点で落とさない要素}
\begin{itemize}
\item 配送例のReal-Worldを具体化
\item 各構成要素の対応付け
\item 決定変数・制約・目的関数を説明
\item 実装とフィードバックを説明
\end{itemize}
\vspace{2mm}\hrule\vspace{2mm}

\clearpage
\SlideHeader{Lecture 2: Information and Decision Modeling}{009}{例: 巡回セールスマン問題}
\SlideImage{9}{../Materials/2026/3DM_02_Information-and-Decision-Modeling_SS26.pdf}
\begin{multicols}{2}\footnotesize\raggedcolumns
\NoteSection{日本語訳}
\begin{itemize}
\item 仮説: 旅行時間は一定である。
\item 情報モデル: 旅行時間行列。
\item 意思決定モデル: 時間的考慮を持たない混合整数計画。
\item 実装すると予想外の結果が生じる。通常、計画より時間がかかる。
\item 仮説を、時間依存旅行時間を捉えるように調整する。
\begin{itemize}
\item 情報モデル: 複数の旅行時間行列と遷移に関する考慮。
\item 意思決定モデル: 時刻を考慮した混合整数計画。決定変数は x\_ij(t) のようになる。
\end{itemize}
\end{itemize}
\NoteSection{解説}
TSPでは、顧客をノード、顧客間移動をアーク、距離または旅行時間を重みとして表す。目的は、全顧客を一度ずつ訪問し、総移動時間または距離を最小にする訪問順序を求めることである。\par

情報モデルとしては、顧客集合、デポ、顧客間旅行時間行列が必要である。意思決定モデルとしては、アークを使うかどうかを表す二値変数 x\_\{ij\}、各顧客へ一度入る・一度出る制約、部分巡回除去制約、総旅行時間最小化の目的関数が必要である。\par

時間依存を入れると、旅行時間は定数 d\_\{ij\} ではなく、出発時刻 t に依存する d\_\{ij\}(t) になる。ある顧客を先に訪問するか後に訪問するかで、後続アークの出発時刻が変わり、旅行時間も変わる。\par
\NoteSection{試験対策ポイント}
\begin{itemize}
\item TSPの基本定式化を、変数・目的関数・制約に分けて説明できる。
\item 時間依存TSPでは、旅行時間行列だけでなく時刻依存関数や複数行列が必要になる。
\end{itemize}
\NoteSection{確認問題と解答}
\begin{enumerate}
\item \textbf{TSPの数理モデルを言葉で説明せよ。}\\ノード集合はデポと顧客、アーク集合は顧客間移動を表す。二値変数 x\_\{ij\}=1 は、i から j へ直接移動することを意味する。目的は sum d\_\{ij\}x\_\{ij\} を最小化することである。制約として、各顧客に一度入る、各顧客から一度出る、複数の小さな巡回ができないようにする部分巡回除去制約が必要である。
\item \textbf{時間依存旅行時間がTSPを難しくする理由を説明せよ。}\\静的TSPではアークのコストは固定である。しかし時間依存TSPでは、iからjへのコストが出発時刻に依存する。ある挿入や訪問順変更により後続の到着時刻が変わり、その後の全アークの旅行時間も変わるため、局所的な評価が全体に波及する。
\end{enumerate}
\end{multicols}

\clearpage
\SlideHeader{Lecture 2: Information and Decision Modeling}{010}{アジェンダ}
\SlideImage{10}{../Materials/2026/3DM_02_Information-and-Decision-Modeling_SS26.pdf}
\begin{multicols}{2}\scriptsize\raggedcolumns
\NoteSection{日本語訳}
\begin{itemize}
\item 情報モデリングと意思決定モデリング。
\item ケーススタディ: 時間依存旅行時間を持つ巡回セールスマン問題。
\end{itemize}
\end{multicols}

\clearpage
\SlideHeader{Lecture 2: Information and Decision Modeling}{011}{動機}
\SlideImage{11}{../Materials/2026/3DM_02_Information-and-Decision-Modeling_SS26.pdf}
\begin{multicols}{2}\footnotesize\raggedcolumns
\NoteSection{日本語訳}
\begin{itemize}
\item サービスルーティング。
\item 高価な走行時間を節約する。
\item 時間的コミットメント
\begin{itemize}
\item 顧客サービス経験を向上させる。
\item 顧客の在宅・利用可能性を確保する。
\end{itemize}
\item どちらも旅行時間の正確な考慮を必要とする。
\end{itemize}
\NoteSection{解説}
時間依存旅行時間では、道路や顧客間の旅行時間が時刻で変わる。朝の通勤時間、昼間、夕方では、同じ道路でも速度が違う。したがって一つの平均旅行時間だけでは現実を十分に表せない。\par

情報モデルとしては、時間帯別旅行時間行列、または連続的な旅行時間関数を使う。例: 7-9時、9-15時、15-18時で別の行列を持つ。\par

重要なのは、時間依存性は情報モデルだけでなく意思決定モデルも変える点である。訪問順序を決めるだけでなく、どの時刻にどのアークを通るかが決定に入る。\par
\NoteSection{試験対策ポイント}
\begin{itemize}
\item 詳細な情報モデルほど現実に近いが、データ量・計算複雑性・解釈可能性にコストがある。
\item 時間依存旅行時間は、FCDなどから推定される。
\end{itemize}
\NoteSection{確認問題と解答}
\begin{enumerate}
\item \textbf{なぜ常に最も詳細な旅行時間モデルを選ばないのか。}\\詳細なモデルは時間帯・道路セグメントごとに多くの観測を必要とする。観測が少ないセルでは平均値が不安定になり、過学習のように現実を誤って表す可能性がある。また、意思決定モデルの変数や制約が増え、計算時間が長くなる。したがって、意思決定に有用な粒度を選ぶ必要がある。
\item \textbf{時間帯別旅行時間行列の長所と短所を述べよ。}\\長所は、朝夕の混雑など平均行列では消える変動を反映できることである。短所は、時間帯ごとに十分な観測が必要であり、境界時刻で旅行時間が不連続になる可能性があること、意思決定モデルが複雑になることである。
\end{enumerate}
\end{multicols}

\clearpage
\SlideHeader{Lecture 2: Information and Decision Modeling}{012}{時間依存旅行時間}
\SlideImage{12}{../Materials/2026/3DM_02_Information-and-Decision-Modeling_SS26.pdf}
\begin{multicols}{2}\footnotesize\raggedcolumns
\NoteSection{日本語訳}
\begin{itemize}
\item 交通は時間と空間により変化する。
\item 道路セグメントごとに旅行時間が異なる。
\item 顧客間の最短経路が異なる。
\item ツアー開始時刻により顧客の訪問順序が異なる。
\end{itemize}
\NoteSection{解説}
時間依存旅行時間では、道路や顧客間の旅行時間が時刻で変わる。朝の通勤時間、昼間、夕方では、同じ道路でも速度が違う。したがって一つの平均旅行時間だけでは現実を十分に表せない。\par

情報モデルとしては、時間帯別旅行時間行列、または連続的な旅行時間関数を使う。例: 7-9時、9-15時、15-18時で別の行列を持つ。\par

重要なのは、時間依存性は情報モデルだけでなく意思決定モデルも変える点である。訪問順序を決めるだけでなく、どの時刻にどのアークを通るかが決定に入る。\par
\NoteSection{試験対策ポイント}
\begin{itemize}
\item 詳細な情報モデルほど現実に近いが、データ量・計算複雑性・解釈可能性にコストがある。
\item 時間依存旅行時間は、FCDなどから推定される。
\end{itemize}
\NoteSection{確認問題と解答}
\begin{enumerate}
\item \textbf{なぜ常に最も詳細な旅行時間モデルを選ばないのか。}\\詳細なモデルは時間帯・道路セグメントごとに多くの観測を必要とする。観測が少ないセルでは平均値が不安定になり、過学習のように現実を誤って表す可能性がある。また、意思決定モデルの変数や制約が増え、計算時間が長くなる。したがって、意思決定に有用な粒度を選ぶ必要がある。
\item \textbf{時間帯別旅行時間行列の長所と短所を述べよ。}\\長所は、朝夕の混雑など平均行列では消える変動を反映できることである。短所は、時間帯ごとに十分な観測が必要であり、境界時刻で旅行時間が不連続になる可能性があること、意思決定モデルが複雑になることである。
\end{enumerate}
\end{multicols}

\clearpage
\SlideHeader{Lecture 2: Information and Decision Modeling}{013}{モデリングの概略}
\SlideImage{13}{../Materials/2026/3DM_02_Information-and-Decision-Modeling_SS26.pdf}
\begin{multicols}{2}\footnotesize\raggedcolumns
\NoteSection{日本語訳}
\begin{itemize}
\item 情報モデル
\begin{itemize}
\item 顧客間、デポを含む時間依存旅行時間。
\item 時間依存旅行時間行列、または複数の旅行時間行列。
\end{itemize}
\item 意思決定モデル
\begin{itemize}
\item 目的: 旅行時間を最小化する。
\item 制約: 全顧客をサービスする。
\item 制約: ルーティング制約、部分巡回除去など。
\item 決定変数: 時刻 t に顧客 i と j の間を移動するか。
\end{itemize}
\end{itemize}
\NoteSection{解説}
時間依存旅行時間では、道路や顧客間の旅行時間が時刻で変わる。朝の通勤時間、昼間、夕方では、同じ道路でも速度が違う。したがって一つの平均旅行時間だけでは現実を十分に表せない。\par

情報モデルとしては、時間帯別旅行時間行列、または連続的な旅行時間関数を使う。例: 7-9時、9-15時、15-18時で別の行列を持つ。\par

重要なのは、時間依存性は情報モデルだけでなく意思決定モデルも変える点である。訪問順序を決めるだけでなく、どの時刻にどのアークを通るかが決定に入る。\par
\NoteSection{試験対策ポイント}
\begin{itemize}
\item 詳細な情報モデルほど現実に近いが、データ量・計算複雑性・解釈可能性にコストがある。
\item 時間依存旅行時間は、FCDなどから推定される。
\end{itemize}
\NoteSection{確認問題と解答}
\begin{enumerate}
\item \textbf{なぜ常に最も詳細な旅行時間モデルを選ばないのか。}\\詳細なモデルは時間帯・道路セグメントごとに多くの観測を必要とする。観測が少ないセルでは平均値が不安定になり、過学習のように現実を誤って表す可能性がある。また、意思決定モデルの変数や制約が増え、計算時間が長くなる。したがって、意思決定に有用な粒度を選ぶ必要がある。
\item \textbf{時間帯別旅行時間行列の長所と短所を述べよ。}\\長所は、朝夕の混雑など平均行列では消える変動を反映できることである。短所は、時間帯ごとに十分な観測が必要であり、境界時刻で旅行時間が不連続になる可能性があること、意思決定モデルが複雑になることである。
\end{enumerate}
\end{multicols}

\clearpage
\SlideHeader{Lecture 2: Information and Decision Modeling}{014}{データ収集: Floating Car Data}
\SlideImage{14}{../Materials/2026/3DM_02_Information-and-Decision-Modeling_SS26.pdf}
\begin{multicols}{2}\footnotesize\raggedcolumns
\NoteSection{日本語訳}
\begin{itemize}
\item 時刻と道路セグメントごとの旅行時間。
\item 正確な平均値を計算するには大量のデータが必要である。
\item Floating Car Data: 大規模な車両群が、位置、方向、場合によって速度を短い時間間隔で送信する。
\end{itemize}
\NoteSection{解説}
Floating Car Data (FCD) は、車両から位置・時刻・場合によって速度を収集したデータである。タクシー、配送車、一般車両のGPSログが例である。\par

FCDはそのまま旅行時間行列ではない。位置を道路セグメントに対応付ける map matching、方向の確認、異常値除去、時間帯ごとの集約が必要である。\par

観測数が少ない道路や時間帯では平均速度が不安定になる。そこで、時間帯を広げる、似た道路をまとめる、空間的・時間的な一般化を行うなどの工夫が必要になる。\par
\NoteSection{試験対策ポイント}
\begin{itemize}
\item FCD -> map matching -> セグメント速度 -> 時間帯別旅行時間 -> 顧客間旅行時間行列、という変換を説明できる。
\item データが最初に来るが、データは必ず加工・検証しなければ情報モデルにならない。
\end{itemize}
\NoteSection{確認問題と解答}
\begin{enumerate}
\item \textbf{FCDから旅行時間行列を作る手順を説明せよ。}\\まず車両の位置・時刻・速度を取得する。次に位置を道路セグメントにマッチングし、方向と時刻を確認する。外れ値や停止などを処理し、セグメントと時間帯ごとに速度または旅行時間を集約する。最後に、顧客間経路を道路セグメント列として計算し、セグメント旅行時間を足し合わせて顧客間旅行時間行列を作る。
\item \textbf{小さな観測数が問題になる理由を説明せよ。}\\観測数が少ないと、たまたま渋滞していた一台や異常なGPS点が平均値に大きく影響する。結果として、情報モデルが現実の典型的な状態ではなく偶然の観測を表してしまう。これを避けるには、時間帯や道路分類で集約したり、外れ値処理や補間を行う。
\end{enumerate}
\end{multicols}

\clearpage
\SlideHeader{Lecture 2: Information and Decision Modeling}{015}{FCDの収集と利用}
\SlideImage{15}{../Materials/2026/3DM_02_Information-and-Decision-Modeling_SS26.pdf}
\begin{multicols}{2}\footnotesize\raggedcolumns
\NoteSection{日本語訳}
\begin{itemize}
\item タクシー車両群。
\item ディスパッチャー。
\item 公共交通管理。
\item プロバイダ。
\item インターネット。
\item GPRS, SMS, RDS-TMC, HTTP による通信。
\item およそ1分に1回の通信。
\end{itemize}
\NoteSection{解説}
Floating Car Data (FCD) は、車両から位置・時刻・場合によって速度を収集したデータである。タクシー、配送車、一般車両のGPSログが例である。\par

FCDはそのまま旅行時間行列ではない。位置を道路セグメントに対応付ける map matching、方向の確認、異常値除去、時間帯ごとの集約が必要である。\par

観測数が少ない道路や時間帯では平均速度が不安定になる。そこで、時間帯を広げる、似た道路をまとめる、空間的・時間的な一般化を行うなどの工夫が必要になる。\par
\NoteSection{試験対策ポイント}
\begin{itemize}
\item FCD -> map matching -> セグメント速度 -> 時間帯別旅行時間 -> 顧客間旅行時間行列、という変換を説明できる。
\item データが最初に来るが、データは必ず加工・検証しなければ情報モデルにならない。
\end{itemize}
\NoteSection{確認問題と解答}
\begin{enumerate}
\item \textbf{FCDから旅行時間行列を作る手順を説明せよ。}\\まず車両の位置・時刻・速度を取得する。次に位置を道路セグメントにマッチングし、方向と時刻を確認する。外れ値や停止などを処理し、セグメントと時間帯ごとに速度または旅行時間を集約する。最後に、顧客間経路を道路セグメント列として計算し、セグメント旅行時間を足し合わせて顧客間旅行時間行列を作る。
\item \textbf{小さな観測数が問題になる理由を説明せよ。}\\観測数が少ないと、たまたま渋滞していた一台や異常なGPS点が平均値に大きく影響する。結果として、情報モデルが現実の典型的な状態ではなく偶然の観測を表してしまう。これを避けるには、時間帯や道路分類で集約したり、外れ値処理や補間を行う。
\end{enumerate}
\end{multicols}

\clearpage
\SlideHeader{Lecture 2: Information and Decision Modeling}{016}{FCD: 速度計算}
\SlideImage{16}{../Materials/2026/3DM_02_Information-and-Decision-Modeling_SS26.pdf}
\begin{multicols}{2}\footnotesize\raggedcolumns
\NoteSection{日本語訳}
\begin{itemize}
\item タクシーデータは、道路、方向、速度を直接与えない。
\item 必要な手順
\begin{itemize}
\item タクシー位置を観測する。
\item 位置を道路セグメントに投影する。
\item このセグメントの観測速度を計算する。
\end{itemize}
\end{itemize}
\NoteSection{解説}
Floating Car Data (FCD) は、車両から位置・時刻・場合によって速度を収集したデータである。タクシー、配送車、一般車両のGPSログが例である。\par

FCDはそのまま旅行時間行列ではない。位置を道路セグメントに対応付ける map matching、方向の確認、異常値除去、時間帯ごとの集約が必要である。\par

観測数が少ない道路や時間帯では平均速度が不安定になる。そこで、時間帯を広げる、似た道路をまとめる、空間的・時間的な一般化を行うなどの工夫が必要になる。\par
\NoteSection{試験対策ポイント}
\begin{itemize}
\item FCD -> map matching -> セグメント速度 -> 時間帯別旅行時間 -> 顧客間旅行時間行列、という変換を説明できる。
\item データが最初に来るが、データは必ず加工・検証しなければ情報モデルにならない。
\end{itemize}
\NoteSection{確認問題と解答}
\begin{enumerate}
\item \textbf{FCDから旅行時間行列を作る手順を説明せよ。}\\まず車両の位置・時刻・速度を取得する。次に位置を道路セグメントにマッチングし、方向と時刻を確認する。外れ値や停止などを処理し、セグメントと時間帯ごとに速度または旅行時間を集約する。最後に、顧客間経路を道路セグメント列として計算し、セグメント旅行時間を足し合わせて顧客間旅行時間行列を作る。
\item \textbf{小さな観測数が問題になる理由を説明せよ。}\\観測数が少ないと、たまたま渋滞していた一台や異常なGPS点が平均値に大きく影響する。結果として、情報モデルが現実の典型的な状態ではなく偶然の観測を表してしまう。これを避けるには、時間帯や道路分類で集約したり、外れ値処理や補間を行う。
\end{enumerate}
\end{multicols}

\clearpage
\SlideHeader{Lecture 2: Information and Decision Modeling}{017}{FCD: 速度計算の表}
\SlideImage{17}{../Materials/2026/3DM_02_Information-and-Decision-Modeling_SS26.pdf}
\begin{multicols}{2}\footnotesize\raggedcolumns
\NoteSection{日本語訳}
\begin{itemize}
\item link, time, speed の観測表。
\item 同じリンクでも時刻により速度が大きく異なる。
\end{itemize}
\NoteSection{解説}
Floating Car Data (FCD) は、車両から位置・時刻・場合によって速度を収集したデータである。タクシー、配送車、一般車両のGPSログが例である。\par

FCDはそのまま旅行時間行列ではない。位置を道路セグメントに対応付ける map matching、方向の確認、異常値除去、時間帯ごとの集約が必要である。\par

観測数が少ない道路や時間帯では平均速度が不安定になる。そこで、時間帯を広げる、似た道路をまとめる、空間的・時間的な一般化を行うなどの工夫が必要になる。\par
\NoteSection{試験対策ポイント}
\begin{itemize}
\item FCD -> map matching -> セグメント速度 -> 時間帯別旅行時間 -> 顧客間旅行時間行列、という変換を説明できる。
\item データが最初に来るが、データは必ず加工・検証しなければ情報モデルにならない。
\end{itemize}
\NoteSection{確認問題と解答}
\begin{enumerate}
\item \textbf{FCDから旅行時間行列を作る手順を説明せよ。}\\まず車両の位置・時刻・速度を取得する。次に位置を道路セグメントにマッチングし、方向と時刻を確認する。外れ値や停止などを処理し、セグメントと時間帯ごとに速度または旅行時間を集約する。最後に、顧客間経路を道路セグメント列として計算し、セグメント旅行時間を足し合わせて顧客間旅行時間行列を作る。
\item \textbf{小さな観測数が問題になる理由を説明せよ。}\\観測数が少ないと、たまたま渋滞していた一台や異常なGPS点が平均値に大きく影響する。結果として、情報モデルが現実の典型的な状態ではなく偶然の観測を表してしまう。これを避けるには、時間帯や道路分類で集約したり、外れ値処理や補間を行う。
\end{enumerate}
\end{multicols}

\clearpage
\ExamHeader{Floating Car Data と旅行時間情報モデル}
\ExamQuestion{WS22/23 Aufgabe 4 [8 Punkte]}
\ExamSubsection{問題内容}
交通観測データから、時間帯別の平均旅行時間または旅行時間行列を構築する手順を説明する問題。\par
\ExamSubsection{満点答案}
まず、車両から時刻付き位置データ、速度、方向などを収集する。次に、位置データを道路ネットワーク上のリンクに対応付ける map matching を行う。これにより、各観測がどの道路セグメント・どの方向・どの時刻に属するかが分かる。\par
その後、エラー、外れ値、停止中データ、GPSのずれを除去または補正する。観測を時間帯と道路セグメントごとにグループ化し、平均速度または平均旅行時間を計算する。観測数が少ないセグメントや時間帯では、時間帯をまとめる、類似道路で補完する、空間的・時間的な一般化を使う。\par
最後に、顧客間の経路を道路セグメント列として求め、各セグメントの時間帯別旅行時間を足し合わせて、顧客間の時間依存旅行時間行列を作る。この情報モデルは、時間依存TSP/VRPの意思決定モデルの入力になる。\par
\ExamSubsection{具体例による解説}
タクシーFCDが1分ごとに位置を送る場合、Hildesheimer Strasseの月曜8時台に属する観測を集め、平均速度を出す。その速度からリンク旅行時間を求め、顧客AからBへ行く経路上のリンク旅行時間を合計してA-Bの8時台旅行時間を作る。\par
\ExamSubsection{採点で落とさない要素}
\begin{itemize}
\item データ収集
\item map matching
\item 外れ値処理
\item 時間帯・セグメント別集約
\item 顧客間行列への変換
\item 観測不足への対応
\end{itemize}
\vspace{2mm}\hrule\vspace{2mm}
\ExamQuestion{SS24 Aufgabe 1b [4 Punkte]}
\ExamSubsection{問題内容}
旅行時間の情報モデルを二つ例示し、なぜ常により詳細なモデルを選ばないのかを説明する問題。\par
\ExamSubsection{満点答案}
旅行時間情報モデルの例として、第一に単一の平均旅行時間行列がある。これは各顧客間の旅行時間を一つの値で表す。第二に時間帯別旅行時間行列がある。これは朝、昼、夕方などの時間帯ごとに異なる旅行時間を持つ。さらに詳細には、道路セグメントごとの連続的な時間依存関数も考えられる。\par
常に詳細なモデルを選ばない理由は、データ量、信頼性、計算複雑性、解釈可能性にコストがあるからである。時間帯や道路を細かく分けるほど、各セルの観測数が減り、平均値が不安定になる。また、意思決定モデルの変数・制約が増え、解くのが難しくなる。したがって、目的に十分な粒度を選ぶ必要がある。\par
\ExamSubsection{具体例による解説}
夜間配送が中心なら、朝夕ラッシュを細かく分けたモデルは不要かもしれない。一方、Same-Day配送で夕方ラッシュに多く走るなら、平均行列では遅延を過小評価するため、時間帯別モデルが有用である。\par
\ExamSubsection{採点で落とさない要素}
\begin{itemize}
\item 二つ以上の情報モデルを提示
\item 詳細モデルの利点
\item データ量・信頼性・計算量の欠点
\item 目的に応じた粒度選択
\end{itemize}
\vspace{2mm}\hrule\vspace{2mm}

\clearpage
\SlideHeader{Lecture 2: Information and Decision Modeling}{018}{データ分析}
\SlideImage{18}{../Materials/2026/3DM_02_Information-and-Decision-Modeling_SS26.pdf}
\begin{multicols}{2}\footnotesize\raggedcolumns
\NoteSection{日本語訳}
\begin{itemize}
\item 一日の中で速度パターンが変化する。これは仮説を確認する。
\item 同じ曜日・同じ時刻でもばらつきがある。これは今回の仮説では無視し、次回扱う。
\item 外れ値やエラーが存在し得る。タクシーデータでは特に注意が必要である。
\end{itemize}
\NoteSection{解説}
Floating Car Data (FCD) は、車両から位置・時刻・場合によって速度を収集したデータである。タクシー、配送車、一般車両のGPSログが例である。\par

FCDはそのまま旅行時間行列ではない。位置を道路セグメントに対応付ける map matching、方向の確認、異常値除去、時間帯ごとの集約が必要である。\par

観測数が少ない道路や時間帯では平均速度が不安定になる。そこで、時間帯を広げる、似た道路をまとめる、空間的・時間的な一般化を行うなどの工夫が必要になる。\par
\NoteSection{試験対策ポイント}
\begin{itemize}
\item FCD -> map matching -> セグメント速度 -> 時間帯別旅行時間 -> 顧客間旅行時間行列、という変換を説明できる。
\item データが最初に来るが、データは必ず加工・検証しなければ情報モデルにならない。
\end{itemize}
\NoteSection{確認問題と解答}
\begin{enumerate}
\item \textbf{FCDから旅行時間行列を作る手順を説明せよ。}\\まず車両の位置・時刻・速度を取得する。次に位置を道路セグメントにマッチングし、方向と時刻を確認する。外れ値や停止などを処理し、セグメントと時間帯ごとに速度または旅行時間を集約する。最後に、顧客間経路を道路セグメント列として計算し、セグメント旅行時間を足し合わせて顧客間旅行時間行列を作る。
\item \textbf{小さな観測数が問題になる理由を説明せよ。}\\観測数が少ないと、たまたま渋滞していた一台や異常なGPS点が平均値に大きく影響する。結果として、情報モデルが現実の典型的な状態ではなく偶然の観測を表してしまう。これを避けるには、時間帯や道路分類で集約したり、外れ値処理や補間を行う。
\end{enumerate}
\end{multicols}

\clearpage
\SlideHeader{Lecture 2: Information and Decision Modeling}{019}{データクリーニング}
\SlideImage{19}{../Materials/2026/3DM_02_Information-and-Decision-Modeling_SS26.pdf}
\begin{multicols}{2}\footnotesize\raggedcolumns
\NoteSection{日本語訳}
\begin{itemize}
\item どの観測が外れ値またはエラーかを判断する。
\item 交通信号、天候など多くの外部要因があるため、客観的評価は難しい。
\item データクリーニング例: 制限速度の1.5倍を超える速度観測をすべてフィルタする。
\end{itemize}
\NoteSection{解説}
Floating Car Data (FCD) は、車両から位置・時刻・場合によって速度を収集したデータである。タクシー、配送車、一般車両のGPSログが例である。\par

FCDはそのまま旅行時間行列ではない。位置を道路セグメントに対応付ける map matching、方向の確認、異常値除去、時間帯ごとの集約が必要である。\par

観測数が少ない道路や時間帯では平均速度が不安定になる。そこで、時間帯を広げる、似た道路をまとめる、空間的・時間的な一般化を行うなどの工夫が必要になる。\par
\NoteSection{試験対策ポイント}
\begin{itemize}
\item FCD -> map matching -> セグメント速度 -> 時間帯別旅行時間 -> 顧客間旅行時間行列、という変換を説明できる。
\item データが最初に来るが、データは必ず加工・検証しなければ情報モデルにならない。
\end{itemize}
\NoteSection{確認問題と解答}
\begin{enumerate}
\item \textbf{FCDから旅行時間行列を作る手順を説明せよ。}\\まず車両の位置・時刻・速度を取得する。次に位置を道路セグメントにマッチングし、方向と時刻を確認する。外れ値や停止などを処理し、セグメントと時間帯ごとに速度または旅行時間を集約する。最後に、顧客間経路を道路セグメント列として計算し、セグメント旅行時間を足し合わせて顧客間旅行時間行列を作る。
\item \textbf{小さな観測数が問題になる理由を説明せよ。}\\観測数が少ないと、たまたま渋滞していた一台や異常なGPS点が平均値に大きく影響する。結果として、情報モデルが現実の典型的な状態ではなく偶然の観測を表してしまう。これを避けるには、時間帯や道路分類で集約したり、外れ値処理や補間を行う。
\end{enumerate}
\end{multicols}

\clearpage
\SlideHeader{Lecture 2: Information and Decision Modeling}{020}{データ可視化}
\SlideImage{20}{../Materials/2026/3DM_02_Information-and-Decision-Modeling_SS26.pdf}
\begin{multicols}{2}\footnotesize\raggedcolumns
\NoteSection{日本語訳}
\begin{itemize}
\item 観測の空間分布を見る。
\item よくカバーされている地域がある。
\item 観測が少ない地域もある。これはタクシーデータに由来する。
\item 時間帯についても同様である。
\item 観測を信頼できるか。
\item 集約と一般化が必要である。
\end{itemize}
\NoteSection{解説}
Floating Car Data (FCD) は、車両から位置・時刻・場合によって速度を収集したデータである。タクシー、配送車、一般車両のGPSログが例である。\par

FCDはそのまま旅行時間行列ではない。位置を道路セグメントに対応付ける map matching、方向の確認、異常値除去、時間帯ごとの集約が必要である。\par

観測数が少ない道路や時間帯では平均速度が不安定になる。そこで、時間帯を広げる、似た道路をまとめる、空間的・時間的な一般化を行うなどの工夫が必要になる。\par
\NoteSection{試験対策ポイント}
\begin{itemize}
\item FCD -> map matching -> セグメント速度 -> 時間帯別旅行時間 -> 顧客間旅行時間行列、という変換を説明できる。
\item データが最初に来るが、データは必ず加工・検証しなければ情報モデルにならない。
\end{itemize}
\NoteSection{確認問題と解答}
\begin{enumerate}
\item \textbf{FCDから旅行時間行列を作る手順を説明せよ。}\\まず車両の位置・時刻・速度を取得する。次に位置を道路セグメントにマッチングし、方向と時刻を確認する。外れ値や停止などを処理し、セグメントと時間帯ごとに速度または旅行時間を集約する。最後に、顧客間経路を道路セグメント列として計算し、セグメント旅行時間を足し合わせて顧客間旅行時間行列を作る。
\item \textbf{小さな観測数が問題になる理由を説明せよ。}\\観測数が少ないと、たまたま渋滞していた一台や異常なGPS点が平均値に大きく影響する。結果として、情報モデルが現実の典型的な状態ではなく偶然の観測を表してしまう。これを避けるには、時間帯や道路分類で集約したり、外れ値処理や補間を行う。
\end{enumerate}
\end{multicols}

\clearpage
\SlideHeader{Lecture 2: Information and Decision Modeling}{021}{集約}
\SlideImage{21}{../Materials/2026/3DM_02_Information-and-Decision-Modeling_SS26.pdf}
\begin{multicols}{2}\footnotesize\raggedcolumns
\NoteSection{日本語訳}
\begin{itemize}
\item 時間ステップ W を定義する。例: 24×7。
\item セグメント l、時間ステップ w ごとに、観測速度から平均旅行速度を計算する。
\item W=1: セグメントごとに一つの値。
\item W=24×7: 一日あたり24個、週7日分の値。
\end{itemize}
\NoteSection{解説}
Floating Car Data (FCD) は、車両から位置・時刻・場合によって速度を収集したデータである。タクシー、配送車、一般車両のGPSログが例である。\par

FCDはそのまま旅行時間行列ではない。位置を道路セグメントに対応付ける map matching、方向の確認、異常値除去、時間帯ごとの集約が必要である。\par

観測数が少ない道路や時間帯では平均速度が不安定になる。そこで、時間帯を広げる、似た道路をまとめる、空間的・時間的な一般化を行うなどの工夫が必要になる。\par
\NoteSection{試験対策ポイント}
\begin{itemize}
\item FCD -> map matching -> セグメント速度 -> 時間帯別旅行時間 -> 顧客間旅行時間行列、という変換を説明できる。
\item データが最初に来るが、データは必ず加工・検証しなければ情報モデルにならない。
\end{itemize}
\NoteSection{確認問題と解答}
\begin{enumerate}
\item \textbf{FCDから旅行時間行列を作る手順を説明せよ。}\\まず車両の位置・時刻・速度を取得する。次に位置を道路セグメントにマッチングし、方向と時刻を確認する。外れ値や停止などを処理し、セグメントと時間帯ごとに速度または旅行時間を集約する。最後に、顧客間経路を道路セグメント列として計算し、セグメント旅行時間を足し合わせて顧客間旅行時間行列を作る。
\item \textbf{小さな観測数が問題になる理由を説明せよ。}\\観測数が少ないと、たまたま渋滞していた一台や異常なGPS点が平均値に大きく影響する。結果として、情報モデルが現実の典型的な状態ではなく偶然の観測を表してしまう。これを避けるには、時間帯や道路分類で集約したり、外れ値処理や補間を行う。
\end{enumerate}
\end{multicols}

\clearpage
\SlideHeader{Lecture 2: Information and Decision Modeling}{022}{一般化}
\SlideImage{22}{../Materials/2026/3DM_02_Information-and-Decision-Modeling_SS26.pdf}
\begin{multicols}{2}\footnotesize\raggedcolumns
\NoteSection{日本語訳}
\begin{itemize}
\item 集約結果は数百万個の値になる。
\item 一部の値は少数の観測にしか基づかない。
\item 意思決定モデルは、アクセスしやすく信頼できる情報モデルを必要とする。
\item データ可視化は、異なる道路セグメントで類似した交通品質パターンを示す。
\item 考え方: 典型的な交通品質に基づいて道路セグメントをグループ化する。
\end{itemize}
\NoteSection{解説}
Floating Car Data (FCD) は、車両から位置・時刻・場合によって速度を収集したデータである。タクシー、配送車、一般車両のGPSログが例である。\par

FCDはそのまま旅行時間行列ではない。位置を道路セグメントに対応付ける map matching、方向の確認、異常値除去、時間帯ごとの集約が必要である。\par

観測数が少ない道路や時間帯では平均速度が不安定になる。そこで、時間帯を広げる、似た道路をまとめる、空間的・時間的な一般化を行うなどの工夫が必要になる。\par
\NoteSection{試験対策ポイント}
\begin{itemize}
\item FCD -> map matching -> セグメント速度 -> 時間帯別旅行時間 -> 顧客間旅行時間行列、という変換を説明できる。
\item データが最初に来るが、データは必ず加工・検証しなければ情報モデルにならない。
\end{itemize}
\NoteSection{確認問題と解答}
\begin{enumerate}
\item \textbf{FCDから旅行時間行列を作る手順を説明せよ。}\\まず車両の位置・時刻・速度を取得する。次に位置を道路セグメントにマッチングし、方向と時刻を確認する。外れ値や停止などを処理し、セグメントと時間帯ごとに速度または旅行時間を集約する。最後に、顧客間経路を道路セグメント列として計算し、セグメント旅行時間を足し合わせて顧客間旅行時間行列を作る。
\item \textbf{小さな観測数が問題になる理由を説明せよ。}\\観測数が少ないと、たまたま渋滞していた一台や異常なGPS点が平均値に大きく影響する。結果として、情報モデルが現実の典型的な状態ではなく偶然の観測を表してしまう。これを避けるには、時間帯や道路分類で集約したり、外れ値処理や補間を行う。
\end{enumerate}
\end{multicols}

\clearpage
\SlideHeader{Lecture 2: Information and Decision Modeling}{023}{一般化: 正規化}
\SlideImage{23}{../Materials/2026/3DM_02_Information-and-Decision-Modeling_SS26.pdf}
\begin{multicols}{2}\footnotesize\raggedcolumns
\NoteSection{日本語訳}
\begin{itemize}
\item 異なるセグメントは同じパターンを持っていても絶対値が異なる。例: 制限速度30または50。
\item 比較を可能にするため、日内最大値で正規化する。
\end{itemize}
\NoteSection{解説}
Floating Car Data (FCD) は、車両から位置・時刻・場合によって速度を収集したデータである。タクシー、配送車、一般車両のGPSログが例である。\par

FCDはそのまま旅行時間行列ではない。位置を道路セグメントに対応付ける map matching、方向の確認、異常値除去、時間帯ごとの集約が必要である。\par

観測数が少ない道路や時間帯では平均速度が不安定になる。そこで、時間帯を広げる、似た道路をまとめる、空間的・時間的な一般化を行うなどの工夫が必要になる。\par
\NoteSection{試験対策ポイント}
\begin{itemize}
\item FCD -> map matching -> セグメント速度 -> 時間帯別旅行時間 -> 顧客間旅行時間行列、という変換を説明できる。
\item データが最初に来るが、データは必ず加工・検証しなければ情報モデルにならない。
\end{itemize}
\NoteSection{確認問題と解答}
\begin{enumerate}
\item \textbf{FCDから旅行時間行列を作る手順を説明せよ。}\\まず車両の位置・時刻・速度を取得する。次に位置を道路セグメントにマッチングし、方向と時刻を確認する。外れ値や停止などを処理し、セグメントと時間帯ごとに速度または旅行時間を集約する。最後に、顧客間経路を道路セグメント列として計算し、セグメント旅行時間を足し合わせて顧客間旅行時間行列を作る。
\item \textbf{小さな観測数が問題になる理由を説明せよ。}\\観測数が少ないと、たまたま渋滞していた一台や異常なGPS点が平均値に大きく影響する。結果として、情報モデルが現実の典型的な状態ではなく偶然の観測を表してしまう。これを避けるには、時間帯や道路分類で集約したり、外れ値処理や補間を行う。
\end{enumerate}
\end{multicols}

\clearpage
\SlideHeader{Lecture 2: Information and Decision Modeling}{024}{一般化: クラスタリング}
\SlideImage{24}{../Materials/2026/3DM_02_Information-and-Decision-Modeling_SS26.pdf}
\begin{multicols}{2}\footnotesize\raggedcolumns
\NoteSection{日本語訳}
\begin{itemize}
\item クラスタリングアルゴリズムで類似した日内変化をグループ化する。
\item 6つのクラスタ。
\item 0は交通品質が悪く旅行時間が長い。
\item 1は交通品質が良く旅行時間が短い。
\end{itemize}
\NoteSection{解説}
Floating Car Data (FCD) は、車両から位置・時刻・場合によって速度を収集したデータである。タクシー、配送車、一般車両のGPSログが例である。\par

FCDはそのまま旅行時間行列ではない。位置を道路セグメントに対応付ける map matching、方向の確認、異常値除去、時間帯ごとの集約が必要である。\par

観測数が少ない道路や時間帯では平均速度が不安定になる。そこで、時間帯を広げる、似た道路をまとめる、空間的・時間的な一般化を行うなどの工夫が必要になる。\par
\NoteSection{試験対策ポイント}
\begin{itemize}
\item FCD -> map matching -> セグメント速度 -> 時間帯別旅行時間 -> 顧客間旅行時間行列、という変換を説明できる。
\item データが最初に来るが、データは必ず加工・検証しなければ情報モデルにならない。
\end{itemize}
\NoteSection{確認問題と解答}
\begin{enumerate}
\item \textbf{FCDから旅行時間行列を作る手順を説明せよ。}\\まず車両の位置・時刻・速度を取得する。次に位置を道路セグメントにマッチングし、方向と時刻を確認する。外れ値や停止などを処理し、セグメントと時間帯ごとに速度または旅行時間を集約する。最後に、顧客間経路を道路セグメント列として計算し、セグメント旅行時間を足し合わせて顧客間旅行時間行列を作る。
\item \textbf{小さな観測数が問題になる理由を説明せよ。}\\観測数が少ないと、たまたま渋滞していた一台や異常なGPS点が平均値に大きく影響する。結果として、情報モデルが現実の典型的な状態ではなく偶然の観測を表してしまう。これを避けるには、時間帯や道路分類で集約したり、外れ値処理や補間を行う。
\end{enumerate}
\end{multicols}

\clearpage
\SlideHeader{Lecture 2: Information and Decision Modeling}{025}{空間的検証}
\SlideImage{25}{../Materials/2026/3DM_02_Information-and-Decision-Modeling_SS26.pdf}
\begin{multicols}{2}\footnotesize\raggedcolumns
\NoteSection{日本語訳}
\begin{itemize}
\item クラスタはピーク時の挙動が異なる。
\item クラスタは異なる道路セグメントタイプを表す。
\begin{itemize}
\item 市中心部
\item 商業地区
\item 高速道路
\item コンベヤーベルト型の道路
\end{itemize}
\end{itemize}
\NoteSection{解説}
Floating Car Data (FCD) は、車両から位置・時刻・場合によって速度を収集したデータである。タクシー、配送車、一般車両のGPSログが例である。\par

FCDはそのまま旅行時間行列ではない。位置を道路セグメントに対応付ける map matching、方向の確認、異常値除去、時間帯ごとの集約が必要である。\par

観測数が少ない道路や時間帯では平均速度が不安定になる。そこで、時間帯を広げる、似た道路をまとめる、空間的・時間的な一般化を行うなどの工夫が必要になる。\par
\NoteSection{試験対策ポイント}
\begin{itemize}
\item FCD -> map matching -> セグメント速度 -> 時間帯別旅行時間 -> 顧客間旅行時間行列、という変換を説明できる。
\item データが最初に来るが、データは必ず加工・検証しなければ情報モデルにならない。
\end{itemize}
\NoteSection{確認問題と解答}
\begin{enumerate}
\item \textbf{FCDから旅行時間行列を作る手順を説明せよ。}\\まず車両の位置・時刻・速度を取得する。次に位置を道路セグメントにマッチングし、方向と時刻を確認する。外れ値や停止などを処理し、セグメントと時間帯ごとに速度または旅行時間を集約する。最後に、顧客間経路を道路セグメント列として計算し、セグメント旅行時間を足し合わせて顧客間旅行時間行列を作る。
\item \textbf{小さな観測数が問題になる理由を説明せよ。}\\観測数が少ないと、たまたま渋滞していた一台や異常なGPS点が平均値に大きく影響する。結果として、情報モデルが現実の典型的な状態ではなく偶然の観測を表してしまう。これを避けるには、時間帯や道路分類で集約したり、外れ値処理や補間を行う。
\end{enumerate}
\end{multicols}

\clearpage
\SlideHeader{Lecture 2: Information and Decision Modeling}{026}{時間的検証}
\SlideImage{26}{../Materials/2026/3DM_02_Information-and-Decision-Modeling_SS26.pdf}
\begin{multicols}{2}\footnotesize\raggedcolumns
\NoteSection{日本語訳}
\begin{itemize}
\item 一日の中の交通品質。Aは高品質、Fは低品質。
\end{itemize}
\NoteSection{解説}
Floating Car Data (FCD) は、車両から位置・時刻・場合によって速度を収集したデータである。タクシー、配送車、一般車両のGPSログが例である。\par

FCDはそのまま旅行時間行列ではない。位置を道路セグメントに対応付ける map matching、方向の確認、異常値除去、時間帯ごとの集約が必要である。\par

観測数が少ない道路や時間帯では平均速度が不安定になる。そこで、時間帯を広げる、似た道路をまとめる、空間的・時間的な一般化を行うなどの工夫が必要になる。\par
\NoteSection{試験対策ポイント}
\begin{itemize}
\item FCD -> map matching -> セグメント速度 -> 時間帯別旅行時間 -> 顧客間旅行時間行列、という変換を説明できる。
\item データが最初に来るが、データは必ず加工・検証しなければ情報モデルにならない。
\end{itemize}
\NoteSection{確認問題と解答}
\begin{enumerate}
\item \textbf{FCDから旅行時間行列を作る手順を説明せよ。}\\まず車両の位置・時刻・速度を取得する。次に位置を道路セグメントにマッチングし、方向と時刻を確認する。外れ値や停止などを処理し、セグメントと時間帯ごとに速度または旅行時間を集約する。最後に、顧客間経路を道路セグメント列として計算し、セグメント旅行時間を足し合わせて顧客間旅行時間行列を作る。
\item \textbf{小さな観測数が問題になる理由を説明せよ。}\\観測数が少ないと、たまたま渋滞していた一台や異常なGPS点が平均値に大きく影響する。結果として、情報モデルが現実の典型的な状態ではなく偶然の観測を表してしまう。これを避けるには、時間帯や道路分類で集約したり、外れ値処理や補間を行う。
\end{enumerate}
\end{multicols}

\clearpage
\SlideHeader{Lecture 2: Information and Decision Modeling}{027}{空間・時間: 深夜}
\SlideImage{27}{../Materials/2026/3DM_02_Information-and-Decision-Modeling_SS26.pdf}
\begin{multicols}{2}\footnotesize\raggedcolumns
\NoteSection{日本語訳}
\begin{itemize}
\item 月曜日3-4時。
\item セグメントの90\%は空いている。
\item 自由流交通。
\end{itemize}
\NoteSection{解説}
Floating Car Data (FCD) は、車両から位置・時刻・場合によって速度を収集したデータである。タクシー、配送車、一般車両のGPSログが例である。\par

FCDはそのまま旅行時間行列ではない。位置を道路セグメントに対応付ける map matching、方向の確認、異常値除去、時間帯ごとの集約が必要である。\par

観測数が少ない道路や時間帯では平均速度が不安定になる。そこで、時間帯を広げる、似た道路をまとめる、空間的・時間的な一般化を行うなどの工夫が必要になる。\par
\NoteSection{試験対策ポイント}
\begin{itemize}
\item FCD -> map matching -> セグメント速度 -> 時間帯別旅行時間 -> 顧客間旅行時間行列、という変換を説明できる。
\item データが最初に来るが、データは必ず加工・検証しなければ情報モデルにならない。
\end{itemize}
\NoteSection{確認問題と解答}
\begin{enumerate}
\item \textbf{FCDから旅行時間行列を作る手順を説明せよ。}\\まず車両の位置・時刻・速度を取得する。次に位置を道路セグメントにマッチングし、方向と時刻を確認する。外れ値や停止などを処理し、セグメントと時間帯ごとに速度または旅行時間を集約する。最後に、顧客間経路を道路セグメント列として計算し、セグメント旅行時間を足し合わせて顧客間旅行時間行列を作る。
\item \textbf{小さな観測数が問題になる理由を説明せよ。}\\観測数が少ないと、たまたま渋滞していた一台や異常なGPS点が平均値に大きく影響する。結果として、情報モデルが現実の典型的な状態ではなく偶然の観測を表してしまう。これを避けるには、時間帯や道路分類で集約したり、外れ値処理や補間を行う。
\end{enumerate}
\end{multicols}

\clearpage
\SlideHeader{Lecture 2: Information and Decision Modeling}{028}{空間・時間: ラッシュアワー}
\SlideImage{28}{../Materials/2026/3DM_02_Information-and-Decision-Modeling_SS26.pdf}
\begin{multicols}{2}\footnotesize\raggedcolumns
\NoteSection{日本語訳}
\begin{itemize}
\item 月曜日17-18時。
\item 30\%が混雑している。
\item ラッシュアワー。
\end{itemize}
\NoteSection{解説}
Floating Car Data (FCD) は、車両から位置・時刻・場合によって速度を収集したデータである。タクシー、配送車、一般車両のGPSログが例である。\par

FCDはそのまま旅行時間行列ではない。位置を道路セグメントに対応付ける map matching、方向の確認、異常値除去、時間帯ごとの集約が必要である。\par

観測数が少ない道路や時間帯では平均速度が不安定になる。そこで、時間帯を広げる、似た道路をまとめる、空間的・時間的な一般化を行うなどの工夫が必要になる。\par
\NoteSection{試験対策ポイント}
\begin{itemize}
\item FCD -> map matching -> セグメント速度 -> 時間帯別旅行時間 -> 顧客間旅行時間行列、という変換を説明できる。
\item データが最初に来るが、データは必ず加工・検証しなければ情報モデルにならない。
\end{itemize}
\NoteSection{確認問題と解答}
\begin{enumerate}
\item \textbf{FCDから旅行時間行列を作る手順を説明せよ。}\\まず車両の位置・時刻・速度を取得する。次に位置を道路セグメントにマッチングし、方向と時刻を確認する。外れ値や停止などを処理し、セグメントと時間帯ごとに速度または旅行時間を集約する。最後に、顧客間経路を道路セグメント列として計算し、セグメント旅行時間を足し合わせて顧客間旅行時間行列を作る。
\item \textbf{小さな観測数が問題になる理由を説明せよ。}\\観測数が少ないと、たまたま渋滞していた一台や異常なGPS点が平均値に大きく影響する。結果として、情報モデルが現実の典型的な状態ではなく偶然の観測を表してしまう。これを避けるには、時間帯や道路分類で集約したり、外れ値処理や補間を行う。
\end{enumerate}
\end{multicols}

\clearpage
\SlideHeader{Lecture 2: Information and Decision Modeling}{029}{情報モデル: 結論}
\SlideImage{29}{../Materials/2026/3DM_02_Information-and-Decision-Modeling_SS26.pdf}
\begin{multicols}{2}\footnotesize\raggedcolumns
\NoteSection{日本語訳}
\begin{itemize}
\item 道路セグメントと各時間について旅行時間行列を持つ。
\item 例: 最初の時間は長い旅行時間、二番目の時間は短い旅行時間、三番目の時間は最長の旅行時間。
\item この情報に基づいてどのようにルーティングするか。
\end{itemize}
\NoteSection{解説}
時間依存旅行時間では、道路や顧客間の旅行時間が時刻で変わる。朝の通勤時間、昼間、夕方では、同じ道路でも速度が違う。したがって一つの平均旅行時間だけでは現実を十分に表せない。\par

情報モデルとしては、時間帯別旅行時間行列、または連続的な旅行時間関数を使う。例: 7-9時、9-15時、15-18時で別の行列を持つ。\par

重要なのは、時間依存性は情報モデルだけでなく意思決定モデルも変える点である。訪問順序を決めるだけでなく、どの時刻にどのアークを通るかが決定に入る。\par
\NoteSection{試験対策ポイント}
\begin{itemize}
\item 詳細な情報モデルほど現実に近いが、データ量・計算複雑性・解釈可能性にコストがある。
\item 時間依存旅行時間は、FCDなどから推定される。
\end{itemize}
\NoteSection{確認問題と解答}
\begin{enumerate}
\item \textbf{なぜ常に最も詳細な旅行時間モデルを選ばないのか。}\\詳細なモデルは時間帯・道路セグメントごとに多くの観測を必要とする。観測が少ないセルでは平均値が不安定になり、過学習のように現実を誤って表す可能性がある。また、意思決定モデルの変数や制約が増え、計算時間が長くなる。したがって、意思決定に有用な粒度を選ぶ必要がある。
\item \textbf{時間帯別旅行時間行列の長所と短所を述べよ。}\\長所は、朝夕の混雑など平均行列では消える変動を反映できることである。短所は、時間帯ごとに十分な観測が必要であり、境界時刻で旅行時間が不連続になる可能性があること、意思決定モデルが複雑になることである。
\end{enumerate}
\end{multicols}

\clearpage
\SlideHeader{Lecture 2: Information and Decision Modeling}{030}{セグメントから経路へ}
\SlideImage{30}{../Materials/2026/3DM_02_Information-and-Decision-Modeling_SS26.pdf}
\begin{multicols}{2}\footnotesize\raggedcolumns
\NoteSection{日本語訳}
\begin{itemize}
\item 顧客間には複数の潜在的経路がある。
\item 最短経路はDijkstraアルゴリズムで計算される。
\item 時間依存性を考慮する必要がある。つまり時間依存最短経路を探す。
\end{itemize}
\NoteSection{解説}
時間依存旅行時間では、道路や顧客間の旅行時間が時刻で変わる。朝の通勤時間、昼間、夕方では、同じ道路でも速度が違う。したがって一つの平均旅行時間だけでは現実を十分に表せない。\par

情報モデルとしては、時間帯別旅行時間行列、または連続的な旅行時間関数を使う。例: 7-9時、9-15時、15-18時で別の行列を持つ。\par

重要なのは、時間依存性は情報モデルだけでなく意思決定モデルも変える点である。訪問順序を決めるだけでなく、どの時刻にどのアークを通るかが決定に入る。\par
\NoteSection{試験対策ポイント}
\begin{itemize}
\item 詳細な情報モデルほど現実に近いが、データ量・計算複雑性・解釈可能性にコストがある。
\item 時間依存旅行時間は、FCDなどから推定される。
\end{itemize}
\NoteSection{確認問題と解答}
\begin{enumerate}
\item \textbf{なぜ常に最も詳細な旅行時間モデルを選ばないのか。}\\詳細なモデルは時間帯・道路セグメントごとに多くの観測を必要とする。観測が少ないセルでは平均値が不安定になり、過学習のように現実を誤って表す可能性がある。また、意思決定モデルの変数や制約が増え、計算時間が長くなる。したがって、意思決定に有用な粒度を選ぶ必要がある。
\item \textbf{時間帯別旅行時間行列の長所と短所を述べよ。}\\長所は、朝夕の混雑など平均行列では消える変動を反映できることである。短所は、時間帯ごとに十分な観測が必要であり、境界時刻で旅行時間が不連続になる可能性があること、意思決定モデルが複雑になることである。
\end{enumerate}
\end{multicols}

\clearpage
\SlideHeader{Lecture 2: Information and Decision Modeling}{031}{セグメントから経路へ}
\SlideImage{31}{../Materials/2026/3DM_02_Information-and-Decision-Modeling_SS26.pdf}
\begin{multicols}{2}\footnotesize\raggedcolumns
\NoteSection{日本語訳}
\begin{itemize}
\item 異なる時刻に出発すると、異なる道路セグメントを通る可能性がある。
\item 時間依存Dijkstraはこれを考慮する。
\end{itemize}
\NoteSection{解説}
時間依存旅行時間では、道路や顧客間の旅行時間が時刻で変わる。朝の通勤時間、昼間、夕方では、同じ道路でも速度が違う。したがって一つの平均旅行時間だけでは現実を十分に表せない。\par

情報モデルとしては、時間帯別旅行時間行列、または連続的な旅行時間関数を使う。例: 7-9時、9-15時、15-18時で別の行列を持つ。\par

重要なのは、時間依存性は情報モデルだけでなく意思決定モデルも変える点である。訪問順序を決めるだけでなく、どの時刻にどのアークを通るかが決定に入る。\par
\NoteSection{試験対策ポイント}
\begin{itemize}
\item 詳細な情報モデルほど現実に近いが、データ量・計算複雑性・解釈可能性にコストがある。
\item 時間依存旅行時間は、FCDなどから推定される。
\end{itemize}
\NoteSection{確認問題と解答}
\begin{enumerate}
\item \textbf{なぜ常に最も詳細な旅行時間モデルを選ばないのか。}\\詳細なモデルは時間帯・道路セグメントごとに多くの観測を必要とする。観測が少ないセルでは平均値が不安定になり、過学習のように現実を誤って表す可能性がある。また、意思決定モデルの変数や制約が増え、計算時間が長くなる。したがって、意思決定に有用な粒度を選ぶ必要がある。
\item \textbf{時間帯別旅行時間行列の長所と短所を述べよ。}\\長所は、朝夕の混雑など平均行列では消える変動を反映できることである。短所は、時間帯ごとに十分な観測が必要であり、境界時刻で旅行時間が不連続になる可能性があること、意思決定モデルが複雑になることである。
\end{enumerate}
\end{multicols}

\clearpage
\SlideHeader{Lecture 2: Information and Decision Modeling}{032}{セグメントから経路へ: FIFO}
\SlideImage{32}{../Materials/2026/3DM_02_Information-and-Decision-Modeling_SS26.pdf}
\begin{multicols}{2}\footnotesize\raggedcolumns
\NoteSection{日本語訳}
\begin{itemize}
\item 行列の硬い切り替えは奇妙な遷移現象を引き起こす。
\item FIFO: 遅く出発したときに追い越しが起きないようにする必要がある。
\item ブレークポイントを線形結合で滑らかにする。
\end{itemize}
\NoteSection{解説}
DDDMでは、データを眺めるだけでなく、データに基づいて実行可能な行動を選ぶ。例えば、食品配送なら『どの注文をどのドライバーに割り当てるか』、在庫管理なら『いつ何個発注するか』が決定である。\par

stochasticity は、将来情報が確実には分からないことを意味する。例として、注文発生、道路混雑、サービス時間、ドライバーのキャンセルがある。dynamism は、新しい情報が入るたびに状態と計画が変わることである。\par

試験答案では、stochasticity=不確実な情報実現、dynamism=時間を通じた状態変化と再意思決定、と分けて書き、最後に『そのため将来を先読みしつつ柔軟性を残す必要がある』と結ぶとよい。\par
\NoteSection{試験対策ポイント}
\begin{itemize}
\item DDDM = データ分析 + 意思決定モデル + 現実への実装、というループで理解する。
\item 単なる予測ではなく、予測・観測を意思決定に変換する点が講義の中心である。
\end{itemize}
\NoteSection{確認問題と解答}
\begin{enumerate}
\item \textbf{DDDMは通常のデータ分析と何が違うか。}\\通常のデータ分析は、過去データの記述や将来予測で終わる場合がある。DDDMでは、その予測や観測を使って具体的な行動を選ぶ。配送なら需要予測だけでなく、車両割当、訪問順序、注文受諾、再計画まで含める。したがって情報モデル、意思決定モデル、実装を一体で説明する必要がある。
\item \textbf{stochasticity と dynamism を例で区別せよ。}\\食品配送で『次にどこから注文が来るか分からない』ことは stochasticity である。注文が実際に入った後、ドライバー位置や未配送注文が変わり、計画を作り直す必要があることは dynamism である。両者が合わさると stochastic dynamic decision problem になる。
\end{enumerate}
\end{multicols}

\clearpage
\SlideHeader{Lecture 2: Information and Decision Modeling}{033}{モデリングの概略}
\SlideImage{33}{../Materials/2026/3DM_02_Information-and-Decision-Modeling_SS26.pdf}
\begin{multicols}{2}\footnotesize\raggedcolumns
\NoteSection{日本語訳}
\begin{itemize}
\item 情報モデル
\begin{itemize}
\item 顧客間、デポを含む時間依存旅行時間。
\item 時間依存旅行時間行列、または複数の旅行時間行列。
\end{itemize}
\item 意思決定モデル
\begin{itemize}
\item 目的: 旅行時間を最小化する。
\item 制約: 全顧客をサービスする。
\item 制約: ルーティング制約、部分巡回除去など。
\item 決定変数: 時刻 t に顧客 i と j の間を移動するか。
\end{itemize}
\end{itemize}
\NoteSection{解説}
時間依存旅行時間では、道路や顧客間の旅行時間が時刻で変わる。朝の通勤時間、昼間、夕方では、同じ道路でも速度が違う。したがって一つの平均旅行時間だけでは現実を十分に表せない。\par

情報モデルとしては、時間帯別旅行時間行列、または連続的な旅行時間関数を使う。例: 7-9時、9-15時、15-18時で別の行列を持つ。\par

重要なのは、時間依存性は情報モデルだけでなく意思決定モデルも変える点である。訪問順序を決めるだけでなく、どの時刻にどのアークを通るかが決定に入る。\par
\NoteSection{試験対策ポイント}
\begin{itemize}
\item 詳細な情報モデルほど現実に近いが、データ量・計算複雑性・解釈可能性にコストがある。
\item 時間依存旅行時間は、FCDなどから推定される。
\end{itemize}
\NoteSection{確認問題と解答}
\begin{enumerate}
\item \textbf{なぜ常に最も詳細な旅行時間モデルを選ばないのか。}\\詳細なモデルは時間帯・道路セグメントごとに多くの観測を必要とする。観測が少ないセルでは平均値が不安定になり、過学習のように現実を誤って表す可能性がある。また、意思決定モデルの変数や制約が増え、計算時間が長くなる。したがって、意思決定に有用な粒度を選ぶ必要がある。
\item \textbf{時間帯別旅行時間行列の長所と短所を述べよ。}\\長所は、朝夕の混雑など平均行列では消える変動を反映できることである。短所は、時間帯ごとに十分な観測が必要であり、境界時刻で旅行時間が不連続になる可能性があること、意思決定モデルが複雑になることである。
\end{enumerate}
\end{multicols}

\clearpage
\SlideHeader{Lecture 2: Information and Decision Modeling}{034}{時間依存決定変数}
\SlideImage{34}{../Materials/2026/3DM_02_Information-and-Decision-Modeling_SS26.pdf}
\begin{multicols}{2}\footnotesize\raggedcolumns
\NoteSection{日本語訳}
\begin{itemize}
\item アークごとの旅行時間は時刻により異なる。
\item 例: d\_AB は時刻により 1, 1, 欠落のような値を取る。
\item 時刻を含むネットワーク表現が必要になる。
\end{itemize}
\NoteSection{解説}
時間依存旅行時間では、道路や顧客間の旅行時間が時刻で変わる。朝の通勤時間、昼間、夕方では、同じ道路でも速度が違う。したがって一つの平均旅行時間だけでは現実を十分に表せない。\par

情報モデルとしては、時間帯別旅行時間行列、または連続的な旅行時間関数を使う。例: 7-9時、9-15時、15-18時で別の行列を持つ。\par

重要なのは、時間依存性は情報モデルだけでなく意思決定モデルも変える点である。訪問順序を決めるだけでなく、どの時刻にどのアークを通るかが決定に入る。\par
\NoteSection{試験対策ポイント}
\begin{itemize}
\item 詳細な情報モデルほど現実に近いが、データ量・計算複雑性・解釈可能性にコストがある。
\item 時間依存旅行時間は、FCDなどから推定される。
\end{itemize}
\NoteSection{確認問題と解答}
\begin{enumerate}
\item \textbf{なぜ常に最も詳細な旅行時間モデルを選ばないのか。}\\詳細なモデルは時間帯・道路セグメントごとに多くの観測を必要とする。観測が少ないセルでは平均値が不安定になり、過学習のように現実を誤って表す可能性がある。また、意思決定モデルの変数や制約が増え、計算時間が長くなる。したがって、意思決定に有用な粒度を選ぶ必要がある。
\item \textbf{時間帯別旅行時間行列の長所と短所を述べよ。}\\長所は、朝夕の混雑など平均行列では消える変動を反映できることである。短所は、時間帯ごとに十分な観測が必要であり、境界時刻で旅行時間が不連続になる可能性があること、意思決定モデルが複雑になることである。
\end{enumerate}
\end{multicols}

\clearpage
\SlideHeader{Lecture 2: Information and Decision Modeling}{035}{時間依存決定変数}
\SlideImage{35}{../Materials/2026/3DM_02_Information-and-Decision-Modeling_SS26.pdf}
\begin{multicols}{2}\footnotesize\raggedcolumns
\NoteSection{日本語訳}
\begin{itemize}
\item A, B, C, D と時刻 t1 から t7 による時間展開ネットワーク。
\item ある時刻のノードから将来時刻のノードへ移動することで、時刻依存旅行時間を表現する。
\end{itemize}
\NoteSection{解説}
時間依存旅行時間では、道路や顧客間の旅行時間が時刻で変わる。朝の通勤時間、昼間、夕方では、同じ道路でも速度が違う。したがって一つの平均旅行時間だけでは現実を十分に表せない。\par

情報モデルとしては、時間帯別旅行時間行列、または連続的な旅行時間関数を使う。例: 7-9時、9-15時、15-18時で別の行列を持つ。\par

重要なのは、時間依存性は情報モデルだけでなく意思決定モデルも変える点である。訪問順序を決めるだけでなく、どの時刻にどのアークを通るかが決定に入る。\par
\NoteSection{試験対策ポイント}
\begin{itemize}
\item 詳細な情報モデルほど現実に近いが、データ量・計算複雑性・解釈可能性にコストがある。
\item 時間依存旅行時間は、FCDなどから推定される。
\end{itemize}
\NoteSection{確認問題と解答}
\begin{enumerate}
\item \textbf{なぜ常に最も詳細な旅行時間モデルを選ばないのか。}\\詳細なモデルは時間帯・道路セグメントごとに多くの観測を必要とする。観測が少ないセルでは平均値が不安定になり、過学習のように現実を誤って表す可能性がある。また、意思決定モデルの変数や制約が増え、計算時間が長くなる。したがって、意思決定に有用な粒度を選ぶ必要がある。
\item \textbf{時間帯別旅行時間行列の長所と短所を述べよ。}\\長所は、朝夕の混雑など平均行列では消える変動を反映できることである。短所は、時間帯ごとに十分な観測が必要であり、境界時刻で旅行時間が不連続になる可能性があること、意思決定モデルが複雑になることである。
\end{enumerate}
\end{multicols}

\clearpage
\SlideHeader{Lecture 2: Information and Decision Modeling}{036}{解を見つける}
\SlideImage{36}{../Materials/2026/3DM_02_Information-and-Decision-Modeling_SS26.pdf}
\begin{multicols}{2}\footnotesize\raggedcolumns
\NoteSection{日本語訳}
\begin{itemize}
\item 最適計算は難しく、ヒューリスティックが必要である。
\item 4つのヒューリスティック
\begin{itemize}
\item LANTIME: メタヒューリスティック
\item Nearest Neighbor
\item Cheapest Insertion
\item Savings
\end{itemize}
\item 平均値に基づくTSPも解く。
\end{itemize}
\NoteSection{解説}
Nearest Neighbor は、現在地から最も近い未訪問顧客を順に選ぶ貪欲法である。簡単だが、後半に悪い選択が残ることがある。\par

Cheapest Insertion は、未訪問顧客を現在のツアーのどこに挿入すると追加コストが最小かを評価する。静的TSPでは使いやすいが、時間依存旅行時間では挿入が後続時刻を変えるため評価が難しくなる。\par

Savings は、別々に訪問する場合と一つのツアーにまとめる場合の節約量を評価する考え方である。ヒューリスティックは高速だが、最適性保証は弱いことを説明できる必要がある。\par
\NoteSection{試験対策ポイント}
\begin{itemize}
\item 主要ヒューリスティックの直感、長所、時間依存で難しくなる理由を説明できる。
\end{itemize}
\NoteSection{確認問題と解答}
\begin{enumerate}
\item \textbf{Cheapest Insertion の考え方を説明せよ。}\\現在のツアーに対して、まだ訪問していない顧客をどの位置に入れると追加旅行時間が最も小さいかを計算し、その最も安い挿入を繰り返す方法である。静的TSPでは追加コストは局所的に計算しやすい。
\item \textbf{時間依存旅行時間で Cheapest Insertion が難しくなる理由を述べよ。}\\一人の顧客を挿入すると、その後の全顧客への到着時刻が変わる。旅行時間が時刻依存なら、後続アークのコストも変わる。したがって、挿入の効果は局所的な追加距離だけでは評価できず、全体の再評価が必要になる。
\end{enumerate}
\end{multicols}

\clearpage
\SlideHeader{Lecture 2: Information and Decision Modeling}{037}{Nearest Neighbor}
\SlideImage{37}{../Materials/2026/3DM_02_Information-and-Decision-Modeling_SS26.pdf}
\begin{multicols}{2}\footnotesize\raggedcolumns
\NoteSection{日本語訳}
\begin{itemize}
\item 現在位置から最も近い顧客を次に選ぶ。
\item 時間依存性を扱うことができる。
\item 通常、終盤に長い移動が残る。
\end{itemize}
\NoteSection{解説}
Nearest Neighbor は、現在地から最も近い未訪問顧客を順に選ぶ貪欲法である。簡単だが、後半に悪い選択が残ることがある。\par

Cheapest Insertion は、未訪問顧客を現在のツアーのどこに挿入すると追加コストが最小かを評価する。静的TSPでは使いやすいが、時間依存旅行時間では挿入が後続時刻を変えるため評価が難しくなる。\par

Savings は、別々に訪問する場合と一つのツアーにまとめる場合の節約量を評価する考え方である。ヒューリスティックは高速だが、最適性保証は弱いことを説明できる必要がある。\par
\NoteSection{試験対策ポイント}
\begin{itemize}
\item 主要ヒューリスティックの直感、長所、時間依存で難しくなる理由を説明できる。
\end{itemize}
\NoteSection{確認問題と解答}
\begin{enumerate}
\item \textbf{Cheapest Insertion の考え方を説明せよ。}\\現在のツアーに対して、まだ訪問していない顧客をどの位置に入れると追加旅行時間が最も小さいかを計算し、その最も安い挿入を繰り返す方法である。静的TSPでは追加コストは局所的に計算しやすい。
\item \textbf{時間依存旅行時間で Cheapest Insertion が難しくなる理由を述べよ。}\\一人の顧客を挿入すると、その後の全顧客への到着時刻が変わる。旅行時間が時刻依存なら、後続アークのコストも変わる。したがって、挿入の効果は局所的な追加距離だけでは評価できず、全体の再評価が必要になる。
\end{enumerate}
\end{multicols}

\clearpage
\SlideHeader{Lecture 2: Information and Decision Modeling}{038}{Cheapest Insertion}
\SlideImage{38}{../Materials/2026/3DM_02_Information-and-Decision-Modeling_SS26.pdf}
\begin{multicols}{2}\footnotesize\raggedcolumns
\NoteSection{日本語訳}
\begin{itemize}
\item 空のツアーから始める。
\item 段階的に最も安い顧客をツアーに挿入する。
\item 顧客がずらされ、旅行時間の変化により以前の決定が有効でなくなる。
\end{itemize}
\NoteSection{解説}
Nearest Neighbor は、現在地から最も近い未訪問顧客を順に選ぶ貪欲法である。簡単だが、後半に悪い選択が残ることがある。\par

Cheapest Insertion は、未訪問顧客を現在のツアーのどこに挿入すると追加コストが最小かを評価する。静的TSPでは使いやすいが、時間依存旅行時間では挿入が後続時刻を変えるため評価が難しくなる。\par

Savings は、別々に訪問する場合と一つのツアーにまとめる場合の節約量を評価する考え方である。ヒューリスティックは高速だが、最適性保証は弱いことを説明できる必要がある。\par
\NoteSection{試験対策ポイント}
\begin{itemize}
\item 主要ヒューリスティックの直感、長所、時間依存で難しくなる理由を説明できる。
\end{itemize}
\NoteSection{確認問題と解答}
\begin{enumerate}
\item \textbf{Cheapest Insertion の考え方を説明せよ。}\\現在のツアーに対して、まだ訪問していない顧客をどの位置に入れると追加旅行時間が最も小さいかを計算し、その最も安い挿入を繰り返す方法である。静的TSPでは追加コストは局所的に計算しやすい。
\item \textbf{時間依存旅行時間で Cheapest Insertion が難しくなる理由を述べよ。}\\一人の顧客を挿入すると、その後の全顧客への到着時刻が変わる。旅行時間が時刻依存なら、後続アークのコストも変わる。したがって、挿入の効果は局所的な追加距離だけでは評価できず、全体の再評価が必要になる。
\end{enumerate}
\end{multicols}

\clearpage
\SlideHeader{Lecture 2: Information and Decision Modeling}{039}{Savings}
\SlideImage{39}{../Materials/2026/3DM_02_Information-and-Decision-Modeling_SS26.pdf}
\begin{multicols}{2}\footnotesize\raggedcolumns
\NoteSection{日本語訳}
\begin{itemize}
\item 直接往復ツアーから始める。
\item 最小費用でツアーを統合する。
\item 時間依存性の評価は難しい。
\item 少なくとも後のステップでは暗黙に考慮される。
\end{itemize}
\NoteSection{解説}
Nearest Neighbor は、現在地から最も近い未訪問顧客を順に選ぶ貪欲法である。簡単だが、後半に悪い選択が残ることがある。\par

Cheapest Insertion は、未訪問顧客を現在のツアーのどこに挿入すると追加コストが最小かを評価する。静的TSPでは使いやすいが、時間依存旅行時間では挿入が後続時刻を変えるため評価が難しくなる。\par

Savings は、別々に訪問する場合と一つのツアーにまとめる場合の節約量を評価する考え方である。ヒューリスティックは高速だが、最適性保証は弱いことを説明できる必要がある。\par
\NoteSection{試験対策ポイント}
\begin{itemize}
\item 主要ヒューリスティックの直感、長所、時間依存で難しくなる理由を説明できる。
\end{itemize}
\NoteSection{確認問題と解答}
\begin{enumerate}
\item \textbf{Cheapest Insertion の考え方を説明せよ。}\\現在のツアーに対して、まだ訪問していない顧客をどの位置に入れると追加旅行時間が最も小さいかを計算し、その最も安い挿入を繰り返す方法である。静的TSPでは追加コストは局所的に計算しやすい。
\item \textbf{時間依存旅行時間で Cheapest Insertion が難しくなる理由を述べよ。}\\一人の顧客を挿入すると、その後の全顧客への到着時刻が変わる。旅行時間が時刻依存なら、後続アークのコストも変わる。したがって、挿入の効果は局所的な追加距離だけでは評価できず、全体の再評価が必要になる。
\end{enumerate}
\end{multicols}

\clearpage
\ExamHeader{TSP、時間依存TSP、ヒューリスティック}
\ExamQuestion{WS22/23 Aufgabe 1 [7 Punkte]}
\ExamSubsection{問題内容}
TSPの数理最適化モデルをスケッチし、時間次元をTSPモデルでどのように考慮するかを説明する問題。\par
\ExamSubsection{満点答案}
標準TSPでは、ノード集合を顧客とデポ、アーク集合を顧客間移動として定義する。各アーク (i,j) には距離または旅行時間 d\_\{ij\} がある。二値決定変数 x\_\{ij\} は、i から j へ直接移動するなら1、そうでなければ0である。目的関数は、sum\_i sum\_j d\_\{ij\} x\_\{ij\} を最小化する。制約として、各顧客にちょうど一度入る、各顧客からちょうど一度出る、部分巡回を排除する制約が必要である。\par
時間次元を考慮する場合、旅行時間 d\_\{ij\} は定数ではなく、出発時刻 t に依存する d\_\{ij\}(t) になる。さらに、顧客 i への到着時刻、サービス開始時刻、顧客 j への出発時刻などの時間変数が必要になる。アークを選ぶかどうかだけでなく、いつそのアークを通るかが重要になるため、x\_\{ij\}(t) のような時刻依存決定変数や時間伝播制約が必要になる。\par
\ExamSubsection{具体例による解説}
A->B が朝8時なら20分、昼12時なら10分なら、同じ訪問順でも出発時刻により目的関数値が変わる。Aを先に訪問するか後に訪問するかが、後続の全旅行時間に影響する。\par
\ExamSubsection{採点で落とさない要素}
\begin{itemize}
\item 変数 x\_ij
\item 目的関数
\item 入次数・出次数制約
\item 部分巡回除去
\item 時間依存旅行時間 d\_ij(t)
\item 到着時刻/出発時刻の伝播
\end{itemize}
\vspace{2mm}\hrule\vspace{2mm}
\ExamQuestion{WS23/24 Aufgabe 5 [7 Punkte]}
\ExamSubsection{問題内容}
Cheapest Insertion が時間依存旅行時間を持つ場合になぜ難しくなるかを説明する問題。\par
\ExamSubsection{満点答案}
Cheapest Insertion は、未訪問顧客を現在のツアーのどの位置へ挿入すると追加コストが最小になるかを計算し、その最安挿入を繰り返すヒューリスティックである。静的TSPでは、追加コストは d\_\{i,k\}+d\_\{k,j\}-d\_\{i,j\} のように局所的に計算できる。\par
時間依存旅行時間では、顧客 k を i と j の間に挿入すると、j 以降の到着時刻が変わる。旅行時間が出発時刻に依存するため、後続アークの旅行時間もすべて変わり得る。したがって、局所的な追加コストだけでは真の影響を測れず、ツアー全体の時刻を再計算しなければならない。\par
\ExamSubsection{具体例による解説}
A-Bの間にCを入れると、Bへの到着が15分遅れ、その結果B-Dを夕方ラッシュ時に走ることになるかもしれない。この場合、Cの挿入自体は短く見えても、後続アークが大きく悪化する。\par
\ExamSubsection{採点で落とさない要素}
\begin{itemize}
\item Cheapest Insertionの基本説明
\item 静的追加コスト
\item 後続到着時刻への波及
\item 時間依存コストの再評価
\item 局所貪欲法の限界
\end{itemize}
\vspace{2mm}\hrule\vspace{2mm}

\clearpage
\SlideHeader{Lecture 2: Information and Decision Modeling}{040}{結果}
\SlideImage{40}{../Materials/2026/3DM_02_Information-and-Decision-Modeling_SS26.pdf}
\begin{multicols}{2}\footnotesize\raggedcolumns
\NoteSection{日本語訳}
\begin{itemize}
\item 出発時刻が異なる。
\item 旅行時間は時刻に依存する。
\item 静的モデルは過大評価も過小評価も行う。
\end{itemize}
\NoteSection{解説}
時間依存旅行時間では、道路や顧客間の旅行時間が時刻で変わる。朝の通勤時間、昼間、夕方では、同じ道路でも速度が違う。したがって一つの平均旅行時間だけでは現実を十分に表せない。\par

情報モデルとしては、時間帯別旅行時間行列、または連続的な旅行時間関数を使う。例: 7-9時、9-15時、15-18時で別の行列を持つ。\par

重要なのは、時間依存性は情報モデルだけでなく意思決定モデルも変える点である。訪問順序を決めるだけでなく、どの時刻にどのアークを通るかが決定に入る。\par
\NoteSection{試験対策ポイント}
\begin{itemize}
\item 詳細な情報モデルほど現実に近いが、データ量・計算複雑性・解釈可能性にコストがある。
\item 時間依存旅行時間は、FCDなどから推定される。
\end{itemize}
\NoteSection{確認問題と解答}
\begin{enumerate}
\item \textbf{なぜ常に最も詳細な旅行時間モデルを選ばないのか。}\\詳細なモデルは時間帯・道路セグメントごとに多くの観測を必要とする。観測が少ないセルでは平均値が不安定になり、過学習のように現実を誤って表す可能性がある。また、意思決定モデルの変数や制約が増え、計算時間が長くなる。したがって、意思決定に有用な粒度を選ぶ必要がある。
\item \textbf{時間帯別旅行時間行列の長所と短所を述べよ。}\\長所は、朝夕の混雑など平均行列では消える変動を反映できることである。短所は、時間帯ごとに十分な観測が必要であり、境界時刻で旅行時間が不連続になる可能性があること、意思決定モデルが複雑になることである。
\end{enumerate}
\end{multicols}

\clearpage
\SlideHeader{Lecture 2: Information and Decision Modeling}{041}{実装}
\SlideImage{41}{../Materials/2026/3DM_02_Information-and-Decision-Modeling_SS26.pdf}
\begin{multicols}{2}\footnotesize\raggedcolumns
\NoteSection{日本語訳}
\begin{itemize}
\item 月曜正午に出発
\begin{itemize}
\item 最初は市中心部を避ける。
\item 先に郊外をサービスする。
\end{itemize}
\item 月曜夕方に出発
\begin{itemize}
\item 先に市中心部をサービスする。
\item 最初に郊外道路を使う。
\item 異なる経路になる。
\end{itemize}
\end{itemize}
\NoteSection{解説}
時間依存旅行時間では、道路や顧客間の旅行時間が時刻で変わる。朝の通勤時間、昼間、夕方では、同じ道路でも速度が違う。したがって一つの平均旅行時間だけでは現実を十分に表せない。\par

情報モデルとしては、時間帯別旅行時間行列、または連続的な旅行時間関数を使う。例: 7-9時、9-15時、15-18時で別の行列を持つ。\par

重要なのは、時間依存性は情報モデルだけでなく意思決定モデルも変える点である。訪問順序を決めるだけでなく、どの時刻にどのアークを通るかが決定に入る。\par
\NoteSection{試験対策ポイント}
\begin{itemize}
\item 詳細な情報モデルほど現実に近いが、データ量・計算複雑性・解釈可能性にコストがある。
\item 時間依存旅行時間は、FCDなどから推定される。
\end{itemize}
\NoteSection{確認問題と解答}
\begin{enumerate}
\item \textbf{なぜ常に最も詳細な旅行時間モデルを選ばないのか。}\\詳細なモデルは時間帯・道路セグメントごとに多くの観測を必要とする。観測が少ないセルでは平均値が不安定になり、過学習のように現実を誤って表す可能性がある。また、意思決定モデルの変数や制約が増え、計算時間が長くなる。したがって、意思決定に有用な粒度を選ぶ必要がある。
\item \textbf{時間帯別旅行時間行列の長所と短所を述べよ。}\\長所は、朝夕の混雑など平均行列では消える変動を反映できることである。短所は、時間帯ごとに十分な観測が必要であり、境界時刻で旅行時間が不連続になる可能性があること、意思決定モデルが複雑になることである。
\end{enumerate}
\end{multicols}

\clearpage
\SlideHeader{Lecture 2: Information and Decision Modeling}{042}{要するに}
\SlideImage{42}{../Materials/2026/3DM_02_Information-and-Decision-Modeling_SS26.pdf}
\begin{multicols}{2}\footnotesize\raggedcolumns
\NoteSection{日本語訳}
\begin{itemize}
\item 問題とデータを真剣に扱うには多くの作業が必要である。
\item 情報モデル: データ収集、分析、構造化、集約、一般化。
\item 意思決定モデル: より多くの決定変数、複雑な意思決定モデル、適応した解法。
\item 得られるアプローチはより詳細なモデリングを行うが、決定論的挙動を仮定している。
\end{itemize}
\NoteSection{解説}
Real-World は、実際の配送現場、道路、顧客、車両、作業員のような現実対象である。現実は非常に複雑なので、そのまま最適化モデルには入れられない。\par

Aggregation は、現実データを意思決定で使える形に圧縮する処理である。例: GPS点列を道路セグメント別速度へ変換する、住所を顧客ノードへ変換する、過去走行記録から旅行時間行列を作る。\par

Information Model は『意思決定に必要な情報の表現』である。顧客集合、デポ、道路ネットワーク、旅行時間行列、需要分布、時間帯などが入る。Decision Model は『その情報を使って何を最適化するか』を表し、決定変数、目的関数、制約、実行可能解を定義する。\par
\NoteSection{試験対策ポイント}
\begin{itemize}
\item Real-World -> Aggregation -> Information Model -> Decision Model -> Implementation -> Real-World の流れを、配送例で説明できる。
\item Information Model と Decision Model を混同しない。前者は情報表現、後者は意思決定の評価・選択構造である。
\end{itemize}
\NoteSection{確認問題と解答}
\begin{enumerate}
\item \textbf{Information Model と Decision Model の違いを説明せよ。}\\Information Model は、現実を意思決定に必要な形へ圧縮した情報表現である。例として、顧客集合、距離行列、時間帯別旅行時間、需要分布がある。Decision Model は、その情報を用いて選択肢を評価する数学的・論理的構造であり、決定変数、目的関数、制約を持つ。配送例なら、旅行時間行列は情報モデル、訪問順序を選んで総旅行時間を最小化するTSP/VRPは意思決定モデルである。
\item \textbf{Aggregation と Implementation を配送例で説明せよ。}\\Aggregation は、GPS記録、道路地図、住所データなどを、顧客ノードや旅行時間行列に変換する処理である。Implementation は、最適化モデルが出した訪問順序やアーク選択を、実際のドライバーが使える配送順、住所リスト、ナビゲーションに変換する処理である。Aggregation は現実からモデルへ、Implementation はモデルから現実への橋である。
\end{enumerate}
\end{multicols}

\clearpage
\SlideHeader{Lecture 2: Information and Decision Modeling}{043}{次回予告}
\SlideImage{43}{../Materials/2026/3DM_02_Information-and-Decision-Modeling_SS26.pdf}
\begin{multicols}{2}\scriptsize\raggedcolumns
\NoteSection{日本語訳}
\begin{itemize}
\item 不確実性と、それが情報モデル・意思決定モデルに与える影響。
\end{itemize}
\end{multicols}
